%% Generated by Sphinx.
\def\sphinxdocclass{report}
\documentclass[a4paper,10pt,ngerman]{sphinxmanual}
\ifdefined\pdfpxdimen
   \let\sphinxpxdimen\pdfpxdimen\else\newdimen\sphinxpxdimen
\fi \sphinxpxdimen=.75bp\relax
\ifdefined\pdfimageresolution
    \pdfimageresolution= \numexpr \dimexpr1in\relax/\sphinxpxdimen\relax
\fi
\newdimen\sphinxremdimen\sphinxremdimen = 10pt
%% let collapsible pdf bookmarks panel have high depth per default
\PassOptionsToPackage{bookmarksdepth=5}{hyperref}

\PassOptionsToPackage{booktabs}{sphinx}
\PassOptionsToPackage{colorrows}{sphinx}

\PassOptionsToPackage{warn}{textcomp}
\usepackage[utf8]{inputenc}
\ifdefined\DeclareUnicodeCharacter
% support both utf8 and utf8x syntaxes
  \ifdefined\DeclareUnicodeCharacterAsOptional
    \def\sphinxDUC#1{\DeclareUnicodeCharacter{"#1}}
  \else
    \let\sphinxDUC\DeclareUnicodeCharacter
  \fi
  \sphinxDUC{00A0}{\nobreakspace}
  \sphinxDUC{2500}{\sphinxunichar{2500}}
  \sphinxDUC{2502}{\sphinxunichar{2502}}
  \sphinxDUC{2514}{\sphinxunichar{2514}}
  \sphinxDUC{251C}{\sphinxunichar{251C}}
  \sphinxDUC{2572}{\textbackslash}
\fi
\usepackage{cmap}
\usepackage[T1]{fontenc}
\usepackage{amsmath,amssymb,amstext}
\usepackage{babel}



\usepackage{tgtermes}
\usepackage{tgheros}
\renewcommand{\ttdefault}{txtt}



\usepackage[Sonny]{fncychap}
\ChNameVar{\Large\normalfont\sffamily}
\ChTitleVar{\Large\normalfont\sffamily}
\usepackage{sphinx}

\fvset{fontsize=auto}
\usepackage{geometry}


% Include hyperref last.
\usepackage{hyperref}
% Fix anchor placement for figures with captions.
\usepackage{hypcap}% it must be loaded after hyperref.
% Set up styles of URL: it should be placed after hyperref.
\urlstyle{same}

\addto\captionsngerman{\renewcommand{\contentsname}{API}}

\usepackage{sphinxmessages}
\setcounter{tocdepth}{1}

\let\cleardoublepage\clearpage

\title{BA\sphinxhyphen{}Sauerteig}
\date{07.07.2026}
\release{1.0}
\author{Patrick Sauerteig}
\newcommand{\sphinxlogo}{\vbox{}}
\renewcommand{\releasename}{Release}
\makeindex
\begin{document}

\ifdefined\shorthandoff
  \ifnum\catcode`\=\string=\active\shorthandoff{=}\fi
  \ifnum\catcode`\"=\active\shorthandoff{"}\fi
\fi

\pagestyle{empty}
\sphinxmaketitle
\pagestyle{plain}
\sphinxtableofcontents
\pagestyle{normal}
\phantomsection\label{\detokenize{index::doc}}


\sphinxstepscope


\chapter{src package}
\label{\detokenize{api:src-package}}\label{\detokenize{api::doc}}

\section{Submodules}
\label{\detokenize{api:submodules}}

\section{src.cost module}
\label{\detokenize{api:module-src.cost}}\label{\detokenize{api:src-cost-module}}\index{module@\spxentry{module}!src.cost@\spxentry{src.cost}}\index{src.cost@\spxentry{src.cost}!module@\spxentry{module}}\index{cost\_function\_whole() (im Modul src.cost)@\spxentry{cost\_function\_whole()}\spxextra{im Modul src.cost}}

\begin{fulllineitems}
\phantomsection\label{\detokenize{api:src.cost.cost_function_whole}}
\pysigstartsignatures
\pysiglinewithargsret
{\sphinxcode{\sphinxupquote{src.cost.}}\sphinxbfcode{\sphinxupquote{cost\_function\_whole}}}
{\sphinxparam{\DUrole{n}{ordering}}\sphinxparamcomma \sphinxparam{\DUrole{n}{node\_grps}}\sphinxparamcomma \sphinxparam{\DUrole{n}{edges}}}
{}
\pysigstopsignatures
\sphinxAtStartPar
Die Kostenfunktion wird in Schritt zwei der Pipeline (Sec. 3.2) verwendet, um die kostengünstigste Achsenordnung zu ermitteln.
Kostengünstig bedeutet hier, dass möglichst viele Achsenpaare mit möglichst geringem Span einfließen sollen.
Diese Funktion wird für Vergleichs\sphinxhyphen{} und Testzwecke weiterhin bereitgestellt.
:param ordering: Ordnung der Achsen (phi)
:type ordering: \DUrole{sphinx_autodoc_typehints-type}{\sphinxcode{\sphinxupquote{list}}{[}\sphinxcode{\sphinxupquote{int}}{]}}
:param node\_grps: key: Achse, value: Knoten auf der Achse
:type node\_grps: \DUrole{sphinx_autodoc_typehints-type}{\sphinxcode{\sphinxupquote{dict}}{[}\sphinxcode{\sphinxupquote{int}}, \sphinxcode{\sphinxupquote{list}}{[}\sphinxcode{\sphinxupquote{int}}, \sphinxcode{\sphinxupquote{str}}{]}{]}}
:param edges: Kantenliste des Graphen
:type edges: \DUrole{sphinx_autodoc_typehints-type}{\sphinxcode{\sphinxupquote{list}}{[}\sphinxcode{\sphinxupquote{tuple}}{[}\sphinxcode{\sphinxupquote{int}} | \sphinxcode{\sphinxupquote{str}}, \sphinxcode{\sphinxupquote{int}} | \sphinxcode{\sphinxupquote{str}}{]}{]}}
\begin{quote}\begin{description}
\sphinxlineitem{Rückgabe}
\sphinxAtStartPar
Gesamtkosten der Achsenordnung

\sphinxlineitem{Rückgabetyp}
\sphinxAtStartPar
\DUrole{sphinx_autodoc_typehints-type}{\sphinxcode{\sphinxupquote{int}}}

\end{description}\end{quote}

\end{fulllineitems}

\index{edges\_between\_axes() (im Modul src.cost)@\spxentry{edges\_between\_axes()}\spxextra{im Modul src.cost}}

\begin{fulllineitems}
\phantomsection\label{\detokenize{api:src.cost.edges_between_axes}}
\pysigstartsignatures
\pysiglinewithargsret
{\sphinxcode{\sphinxupquote{src.cost.}}\sphinxbfcode{\sphinxupquote{edges\_between\_axes}}}
{\sphinxparam{\DUrole{n}{node\_grps}}\sphinxparamcomma \sphinxparam{\DUrole{n}{edges}}\sphinxparamcomma \sphinxparam{\DUrole{n}{a1}}\sphinxparamcomma \sphinxparam{\DUrole{n}{a2}}}
{}
\pysigstopsignatures
\sphinxAtStartPar
Ermittelt wieviele Kanten zwischen den gegebenen Achsen bestehen.
\begin{quote}\begin{description}
\sphinxlineitem{Parameter}\begin{itemize}
\item {} 
\sphinxAtStartPar
\sphinxstyleliteralstrong{\sphinxupquote{node\_grps}} (\DUrole{sphinx_autodoc_typehints-type}{\sphinxcode{\sphinxupquote{dict}}{[}\sphinxcode{\sphinxupquote{int}}, \sphinxcode{\sphinxupquote{list}}{[}\sphinxcode{\sphinxupquote{int}}, \sphinxcode{\sphinxupquote{str}}{]}{]}}) \textendash{} key: subset, value: list

\item {} 
\sphinxAtStartPar
\sphinxstyleliteralstrong{\sphinxupquote{edges}} (\DUrole{sphinx_autodoc_typehints-type}{\sphinxcode{\sphinxupquote{list}}{[}\sphinxcode{\sphinxupquote{tuple}}{[}\sphinxcode{\sphinxupquote{int}} | \sphinxcode{\sphinxupquote{str}}, \sphinxcode{\sphinxupquote{int}} | \sphinxcode{\sphinxupquote{str}}{]}{]}}) \textendash{} Kantenliste aus Tupeln

\item {} 
\sphinxAtStartPar
\sphinxstyleliteralstrong{\sphinxupquote{a1}} (\DUrole{sphinx_autodoc_typehints-type}{\sphinxcode{\sphinxupquote{int}}}) \textendash{} Achse 1

\item {} 
\sphinxAtStartPar
\sphinxstyleliteralstrong{\sphinxupquote{a2}} (\DUrole{sphinx_autodoc_typehints-type}{\sphinxcode{\sphinxupquote{int}}}) \textendash{} Achse 2

\end{itemize}

\sphinxlineitem{Rückgabe}
\sphinxAtStartPar
Anzahl der Kanten

\sphinxlineitem{Rückgabetyp}
\sphinxAtStartPar
\DUrole{sphinx_autodoc_typehints-type}{\sphinxcode{\sphinxupquote{int}}}

\end{description}\end{quote}

\end{fulllineitems}

\index{node\_or\_axes\_span() (im Modul src.cost)@\spxentry{node\_or\_axes\_span()}\spxextra{im Modul src.cost}}

\begin{fulllineitems}
\phantomsection\label{\detokenize{api:src.cost.node_or_axes_span}}
\pysigstartsignatures
\pysiglinewithargsret
{\sphinxcode{\sphinxupquote{src.cost.}}\sphinxbfcode{\sphinxupquote{node\_or\_axes\_span}}}
{\sphinxparam{\DUrole{n}{n1}}\sphinxparamcomma \sphinxparam{\DUrole{n}{n2}}\sphinxparamcomma \sphinxparam{\DUrole{n}{k}}}
{}
\pysigstopsignatures
\sphinxAtStartPar
Berechnet den Span zweier Knoten oder Achsen.
\begin{quote}\begin{description}
\sphinxlineitem{Parameter}\begin{itemize}
\item {} 
\sphinxAtStartPar
\sphinxstyleliteralstrong{\sphinxupquote{n1}} (\DUrole{sphinx_autodoc_typehints-type}{\sphinxcode{\sphinxupquote{int}}}) \textendash{} Achse/Knoten 1

\item {} 
\sphinxAtStartPar
\sphinxstyleliteralstrong{\sphinxupquote{n2}} (\DUrole{sphinx_autodoc_typehints-type}{\sphinxcode{\sphinxupquote{int}}}) \textendash{} Achse/Knoten 2

\item {} 
\sphinxAtStartPar
\sphinxstyleliteralstrong{\sphinxupquote{k}} (\DUrole{sphinx_autodoc_typehints-type}{\sphinxcode{\sphinxupquote{int}}}) \textendash{} Gesamtzahl der Achsen

\end{itemize}

\sphinxlineitem{Rückgabe}
\sphinxAtStartPar
errechneter Spann

\sphinxlineitem{Rückgabetyp}
\sphinxAtStartPar
\DUrole{sphinx_autodoc_typehints-type}{\sphinxcode{\sphinxupquote{int}}}

\end{description}\end{quote}

\end{fulllineitems}



\section{src.crossing\_minimization module}
\label{\detokenize{api:module-src.crossing_minimization}}\label{\detokenize{api:src-crossing-minimization-module}}\index{module@\spxentry{module}!src.crossing\_minimization@\spxentry{src.crossing\_minimization}}\index{src.crossing\_minimization@\spxentry{src.crossing\_minimization}!module@\spxentry{module}}\index{barycenter\_crossmin\_pipeline() (im Modul src.crossing\_minimization)@\spxentry{barycenter\_crossmin\_pipeline()}\spxextra{im Modul src.crossing\_minimization}}

\begin{fulllineitems}
\phantomsection\label{\detokenize{api:src.crossing_minimization.barycenter_crossmin_pipeline}}
\pysigstartsignatures
\pysiglinewithargsret
{\sphinxcode{\sphinxupquote{src.crossing\_minimization.}}\sphinxbfcode{\sphinxupquote{barycenter\_crossmin\_pipeline}}}
{\sphinxparam{\DUrole{n}{layout}}\sphinxparamcomma \sphinxparam{\DUrole{n}{logger}}\sphinxparamcomma \sphinxparam{\DUrole{n}{threshold}\DUrole{o}{=}\DUrole{default_value}{inf}}\sphinxparamcomma \sphinxparam{\DUrole{n}{paper\_like}\DUrole{o}{=}\DUrole{default_value}{True}}}
{}
\pysigstopsignatures
\sphinxAtStartPar
Führt die Kreuzungsminimierungs\sphinxhyphen{}Pipeline mit der Barycenterheuristik aus (3a/3b).

\sphinxAtStartPar
Die Pipeline umfasst im paper\_like = false:
1. anfängliches Layout expandieren
2. Entfernen isolierter Knoten auf den Achsen.
3. Einfügen von Dummy\sphinxhyphen{}Knoten für lange Kanten (Span \textgreater{} 1) und sortieren der inter axis Knoten (Span = 0).
4. Initialisierung des Sweeps und Updates der Datenstrukturen.
5. Mehrfache CW/CCW\sphinxhyphen{}Sweeps für reale und anschließend virtuelle Knoten, die die Barycenterpositionen ermitteln und die Achsen entsprechend sortieren.
6. Herstellen der Achsenornung und letztes Update der Datenstrukturen um die Visualisierung starten zu können.

\sphinxAtStartPar
Die Pipeline umfasst im nicht paper\_like = true:
1. Entfernen isolierter Knoten auf den Achsen.
2. Einfügen von Dummy\sphinxhyphen{}Knoten für lange Kanten (Span \textgreater{} 1) und sortieren der inter axis Knoten (Span = 0).
3. Initialisierung des Sweeps und Updates der Datenstrukturen.
4. Mehrfache CW/CCW\sphinxhyphen{}Sweeps für reale und anschließend virtuelle Knoten, die die Barycenterpositionen ermitteln und die Achsen entsprechend sortieren.
5. Barycenter 3b
6. Herstellen der Achsenornung und letztes Update der Datenstrukturen um die Visualisierung starten zu können.
\begin{quote}\begin{description}
\sphinxlineitem{Parameter}\begin{itemize}
\item {} 
\sphinxAtStartPar
\sphinxstyleliteralstrong{\sphinxupquote{layout}} (\DUrole{sphinx_autodoc_typehints-type}{\sphinxcode{\sphinxupquote{HivePlotLayout}}}) \textendash{} Das zu optimierende HivePlotLayout.

\item {} 
\sphinxAtStartPar
\sphinxstyleliteralstrong{\sphinxupquote{threshold}} (\DUrole{sphinx_autodoc_typehints-type}{\sphinxcode{\sphinxupquote{float}}}) \textendash{} Optionaler Abbruchschwellwert für die Anzahl der Sweep\sphinxhyphen{}Durchläufe.

\item {} 
\sphinxAtStartPar
\sphinxstyleliteralstrong{\sphinxupquote{layout\_expanded}} (\sphinxstyleliteralemphasis{\sphinxupquote{bool}}) \textendash{} dient der Unterscheidung, ob in der Pipeline mit expandierten Achsen gerechnet wird oder nicht, Default = False (nicht expandierter Fall)

\item {} 
\sphinxAtStartPar
\sphinxstyleliteralstrong{\sphinxupquote{paper\_like}} (\DUrole{sphinx_autodoc_typehints-type}{\sphinxcode{\sphinxupquote{bool}}}) \textendash{} dient der Unterscheidung zwischen identischer und unterschiedlicher Knotenordnung

\end{itemize}

\sphinxlineitem{Rückgabetyp}
\sphinxAtStartPar
\DUrole{sphinx_autodoc_typehints-type}{\sphinxcode{\sphinxupquote{None}}}

\end{description}\end{quote}

\end{fulllineitems}

\index{barycenter\_heuristic() (im Modul src.crossing\_minimization)@\spxentry{barycenter\_heuristic()}\spxextra{im Modul src.crossing\_minimization}}

\begin{fulllineitems}
\phantomsection\label{\detokenize{api:src.crossing_minimization.barycenter_heuristic}}
\pysigstartsignatures
\pysiglinewithargsret
{\sphinxcode{\sphinxupquote{src.crossing\_minimization.}}\sphinxbfcode{\sphinxupquote{barycenter\_heuristic}}}
{\sphinxparam{\DUrole{n}{layout}}\sphinxparamcomma \sphinxparam{\DUrole{n}{neighborhood\_map}}\sphinxparamcomma \sphinxparam{\DUrole{n}{node\_axis\_map}}\sphinxparamcomma \sphinxparam{\DUrole{n}{threshold}\DUrole{o}{=}\DUrole{default_value}{inf}}\sphinxparamcomma \sphinxparam{\DUrole{n}{real}\DUrole{o}{=}\DUrole{default_value}{True}}\sphinxparamcomma \sphinxparam{\DUrole{n}{layout\_expanded}\DUrole{o}{=}\DUrole{default_value}{False}}\sphinxparamcomma \sphinxparam{\DUrole{n}{use\_fixed\_positions}\DUrole{o}{=}\DUrole{default_value}{False}}}
{}
\pysigstopsignatures
\sphinxAtStartPar
Führt wiederholte Barycenter\sphinxhyphen{}Sweeps im und gegen den Uhrzeigersinn über alle Achsen aus. Abhängig vom Parameter real werden entweder die realen Knoten (node\_groups)
oder die Dummy\sphinxhyphen{}Knoten (node\_groups\_dummies) pro Achse entsprechend der Barycenterposition ihrer Nachbarn umsortiert. Der Sweep endet, wenn keine Änderung mehr auftritt oder der threshold an Iterationen erreicht ist. Es kann zu Osszilationen im Sweep kommen, wenn
z.B. wenn zwei Knoten die gleiche Barycenter Position haben oder in den Sweeps einfach nur ihre Positionen hin\sphinxhyphen{} und hertauschen. Dies führt dazu, dass die Schleife nie terminiert. In state\_set werden dementsprechend alle erreichten Zustände gehasht und beim Wiederkehren eines zuvor errechneten Zustands kann die Schleife vor einem neuen Durchgang abbrechen.
\begin{quote}\begin{description}
\sphinxlineitem{Parameter}\begin{itemize}
\item {} 
\sphinxAtStartPar
\sphinxstyleliteralstrong{\sphinxupquote{layout}} (\DUrole{sphinx_autodoc_typehints-type}{\sphinxcode{\sphinxupquote{HivePlotLayout}}}) \textendash{} Aktuelles HivePlotLayout mit Achsenzuordnung und Knotenlisten.

\item {} 
\sphinxAtStartPar
\sphinxstyleliteralstrong{\sphinxupquote{neighborhood\_map}} (\DUrole{sphinx_autodoc_typehints-type}{\sphinxcode{\sphinxupquote{dict}}{[}\sphinxcode{\sphinxupquote{int}} | \sphinxcode{\sphinxupquote{str}}, \sphinxcode{\sphinxupquote{list}}{[}\sphinxcode{\sphinxupquote{int}} | \sphinxcode{\sphinxupquote{str}}{]}{]}}) \textendash{} Map von Knoten (real und/oder Dummy, initialisierungsabhängig) auf deren Nachbarn.

\item {} 
\sphinxAtStartPar
\sphinxstyleliteralstrong{\sphinxupquote{node\_axis\_map}} (\DUrole{sphinx_autodoc_typehints-type}{\sphinxcode{\sphinxupquote{dict}}{[}\sphinxcode{\sphinxupquote{int}} | \sphinxcode{\sphinxupquote{str}}, \sphinxcode{\sphinxupquote{int}}{]}}) \textendash{} Map von Knoten auf Achsen in phi.

\item {} 
\sphinxAtStartPar
\sphinxstyleliteralstrong{\sphinxupquote{threshold}} \textendash{} Maximale Anzahl Sweep\sphinxhyphen{}Durchläufe, bevor abgebrochen wird.

\item {} 
\sphinxAtStartPar
\sphinxstyleliteralstrong{\sphinxupquote{real}} (\DUrole{sphinx_autodoc_typehints-type}{\sphinxcode{\sphinxupquote{bool}}}) \textendash{} True, um reale Knoten zu sortieren; False, um Dummy\sphinxhyphen{}Knoten zu sortieren.

\item {} 
\sphinxAtStartPar
\sphinxstyleliteralstrong{\sphinxupquote{layout\_expanded}} (\DUrole{sphinx_autodoc_typehints-type}{\sphinxcode{\sphinxupquote{bool}}}) \textendash{} dient der Unterscheidung, ob in der Pipeline mit expandierten Achsen gerechnet wird oder nicht, Default = False (nicht expandierter Fall)

\item {} 
\sphinxAtStartPar
\sphinxstyleliteralstrong{\sphinxupquote{use\_fixed\_positions}} (\DUrole{sphinx_autodoc_typehints-type}{\sphinxcode{\sphinxupquote{bool}}}) \textendash{} wird für 3b benötigt, damit fixierte positionen verwendet werden

\end{itemize}

\sphinxlineitem{Rückgabetyp}
\sphinxAtStartPar
\DUrole{sphinx_autodoc_typehints-type}{\sphinxcode{\sphinxupquote{None}}}

\end{description}\end{quote}

\end{fulllineitems}

\index{calculate\_barycenter\_position() (im Modul src.crossing\_minimization)@\spxentry{calculate\_barycenter\_position()}\spxextra{im Modul src.crossing\_minimization}}

\begin{fulllineitems}
\phantomsection\label{\detokenize{api:src.crossing_minimization.calculate_barycenter_position}}
\pysigstartsignatures
\pysiglinewithargsret
{\sphinxcode{\sphinxupquote{src.crossing\_minimization.}}\sphinxbfcode{\sphinxupquote{calculate\_barycenter\_position}}}
{\sphinxparam{\DUrole{n}{layout}}\sphinxparamcomma \sphinxparam{\DUrole{n}{neighbor\_group}}\sphinxparamcomma \sphinxparam{\DUrole{n}{node\_axis\_map}}\sphinxparamcomma \sphinxparam{\DUrole{n}{layout\_expanded}\DUrole{o}{=}\DUrole{default_value}{False}}}
{}
\pysigstopsignatures
\sphinxAtStartPar
Berechnet die Barycenterposition eines Knotens über seine ermittelten Nachbarn. Siehe HivePlotLayout.get\_proper\_neighbors().
\begin{quote}\begin{description}
\sphinxlineitem{Parameter}\begin{itemize}
\item {} 
\sphinxAtStartPar
\sphinxstyleliteralstrong{\sphinxupquote{layout}} (\DUrole{sphinx_autodoc_typehints-type}{\sphinxcode{\sphinxupquote{HivePlotLayout}}}) \textendash{} das zu berechnende HivePlotLayout

\item {} 
\sphinxAtStartPar
\sphinxstyleliteralstrong{\sphinxupquote{neighbor\_group}} (\DUrole{sphinx_autodoc_typehints-type}{\sphinxcode{\sphinxupquote{list}}{[}\sphinxcode{\sphinxupquote{int}}{]}}) \textendash{} Liste der Nachbarknoten

\item {} 
\sphinxAtStartPar
\sphinxstyleliteralstrong{\sphinxupquote{layout\_expanded}} (\DUrole{sphinx_autodoc_typehints-type}{\sphinxcode{\sphinxupquote{bool}}}) \textendash{} dient der Unterscheidung, ob in der Pipeline mit expandierten Achsen gerechnet wird oder nicht, Default = False (nicht expandierter Fall)

\end{itemize}

\sphinxlineitem{Rückgabe}
\sphinxAtStartPar
Barycenterposition

\sphinxlineitem{Rückgabetyp}
\sphinxAtStartPar
\DUrole{sphinx_autodoc_typehints-type}{\sphinxcode{\sphinxupquote{float}}}

\end{description}\end{quote}

\end{fulllineitems}

\index{edge\_node\_cleanup() (im Modul src.crossing\_minimization)@\spxentry{edge\_node\_cleanup()}\spxextra{im Modul src.crossing\_minimization}}

\begin{fulllineitems}
\phantomsection\label{\detokenize{api:src.crossing_minimization.edge_node_cleanup}}
\pysigstartsignatures
\pysiglinewithargsret
{\sphinxcode{\sphinxupquote{src.crossing\_minimization.}}\sphinxbfcode{\sphinxupquote{edge\_node\_cleanup}}}
{\sphinxparam{\DUrole{n}{layout}}}
{}
\pysigstopsignatures
\sphinxAtStartPar
Nachbereitung der Pipeline, dient als Absicherung, dass Kanten und intra Knoten im Layout korrekt gesetzt sind.

\end{fulllineitems}

\index{finish\_structured\_axis\_orders() (im Modul src.crossing\_minimization)@\spxentry{finish\_structured\_axis\_orders()}\spxextra{im Modul src.crossing\_minimization}}

\begin{fulllineitems}
\phantomsection\label{\detokenize{api:src.crossing_minimization.finish_structured_axis_orders}}
\pysigstartsignatures
\pysiglinewithargsret
{\sphinxcode{\sphinxupquote{src.crossing\_minimization.}}\sphinxbfcode{\sphinxupquote{finish\_structured\_axis\_orders}}}
{\sphinxparam{\DUrole{n}{layout}}\sphinxparamcomma \sphinxparam{\DUrole{n}{isolated\_node\_groups}}\sphinxparamcomma \sphinxparam{\DUrole{n}{layout\_expanded}\DUrole{o}{=}\DUrole{default_value}{False}}}
{}
\pysigstopsignatures
\sphinxAtStartPar
Aktualisiert die node\_groups am Ende der Pipeline und fügt isolierte Knoten wieder ein.
\begin{quote}\begin{description}
\sphinxlineitem{Rückgabetyp}
\sphinxAtStartPar
\DUrole{sphinx_autodoc_typehints-type}{\sphinxcode{\sphinxupquote{None}}}

\end{description}\end{quote}

\end{fulllineitems}

\index{parse\_dummy\_name() (im Modul src.crossing\_minimization)@\spxentry{parse\_dummy\_name()}\spxextra{im Modul src.crossing\_minimization}}

\begin{fulllineitems}
\phantomsection\label{\detokenize{api:src.crossing_minimization.parse_dummy_name}}
\pysigstartsignatures
\pysiglinewithargsret
{\sphinxcode{\sphinxupquote{src.crossing\_minimization.}}\sphinxbfcode{\sphinxupquote{parse\_dummy\_name}}}
{\sphinxparam{\DUrole{n}{name}}}
{}
\pysigstopsignatures
\sphinxAtStartPar
Zerlegt den Namen eines Dummyknotens in seine Bestandteile: Startknoten, Endknoten und Sequenznummer. (siehe make\_dummy\_name Funktion für Namensschema)
\begin{quote}\begin{description}
\sphinxlineitem{Parameter}
\sphinxAtStartPar
\sphinxstyleliteralstrong{\sphinxupquote{name}} (\DUrole{sphinx_autodoc_typehints-type}{\sphinxcode{\sphinxupquote{str}}}) \textendash{} Der zu zerlegende Dummyknoten

\sphinxlineitem{Rückgabe}
\sphinxAtStartPar
Tupel aus int\sphinxhyphen{}Werten: (Startknoten, Endknoten, Sequenznummer)

\sphinxlineitem{Rückgabetyp}
\sphinxAtStartPar
\DUrole{sphinx_autodoc_typehints-type}{\sphinxcode{\sphinxupquote{tuple}}{[}\sphinxcode{\sphinxupquote{int}}, \sphinxcode{\sphinxupquote{int}}, \sphinxcode{\sphinxupquote{int}}{]}}

\end{description}\end{quote}

\end{fulllineitems}

\index{remove\_isolated\_nodes() (im Modul src.crossing\_minimization)@\spxentry{remove\_isolated\_nodes()}\spxextra{im Modul src.crossing\_minimization}}

\begin{fulllineitems}
\phantomsection\label{\detokenize{api:src.crossing_minimization.remove_isolated_nodes}}
\pysigstartsignatures
\pysiglinewithargsret
{\sphinxcode{\sphinxupquote{src.crossing\_minimization.}}\sphinxbfcode{\sphinxupquote{remove\_isolated\_nodes}}}
{\sphinxparam{\DUrole{n}{graph}}\sphinxparamcomma \sphinxparam{\DUrole{n}{node\_groups}}}
{}
\pysigstopsignatures
\sphinxAtStartPar
Die Funktion entfernt alle isolierten Knoten aus der persistenten node\_group des HivePlotLayouts.
\begin{quote}\begin{description}
\sphinxlineitem{Parameter}\begin{itemize}
\item {} 
\sphinxAtStartPar
\sphinxstyleliteralstrong{\sphinxupquote{graph}} (\DUrole{sphinx_autodoc_typehints-type}{\sphinxcode{\sphinxupquote{Graph}}}) \textendash{} Der zugrundeliegende Graph aus dem HivePlotLayout.

\item {} 
\sphinxAtStartPar
\sphinxstyleliteralstrong{\sphinxupquote{node\_groups}} (\DUrole{sphinx_autodoc_typehints-type}{\sphinxcode{\sphinxupquote{dict}}{[}\sphinxcode{\sphinxupquote{int}}, \sphinxcode{\sphinxupquote{list}}{[}\sphinxcode{\sphinxupquote{int}}{]}{]}}) \textendash{} Die persistente Knotenlisten der Achsen.

\end{itemize}

\sphinxlineitem{Rückgabe}
\sphinxAtStartPar
Ein dict mit den isolierten Knoten pro Achse, key: Achse, value: Listen mit isolierten Knoten.

\sphinxlineitem{Rückgabetyp}
\sphinxAtStartPar
\DUrole{sphinx_autodoc_typehints-type}{\sphinxcode{\sphinxupquote{dict}}{[}\sphinxcode{\sphinxupquote{int}}, \sphinxcode{\sphinxupquote{list}}{[}\sphinxcode{\sphinxupquote{int}}{]}{]}}

\end{description}\end{quote}

\end{fulllineitems}

\index{subdivide\_long\_edges() (im Modul src.crossing\_minimization)@\spxentry{subdivide\_long\_edges()}\spxextra{im Modul src.crossing\_minimization}}

\begin{fulllineitems}
\phantomsection\label{\detokenize{api:src.crossing_minimization.subdivide_long_edges}}
\pysigstartsignatures
\pysiglinewithargsret
{\sphinxcode{\sphinxupquote{src.crossing\_minimization.}}\sphinxbfcode{\sphinxupquote{subdivide\_long\_edges}}}
{\sphinxparam{\DUrole{n}{layout}}\sphinxparamcomma \sphinxparam{\DUrole{n}{node\_position\_map}}\sphinxparamcomma \sphinxparam{\DUrole{n}{node\_axis\_map}}\sphinxparamcomma \sphinxparam{\DUrole{n}{neighborhood\_map}}}
{}
\pysigstopsignatures
\sphinxAtStartPar
Funktion dient dem Einfügen von Dummyknoten für lange Kanten (span \textgreater{} 1) auf den Achsen zwischen Start\sphinxhyphen{} und Endknoten. Schema: d\_{[}Startknoten{]}\_{[}Endknoten{]}\_{[}Sequenznummer{]}: z.B. d\_5\_10\_2 ist der zweite Dummyknoten auf der zerlegten langen Kante von 5 nach 10. Die lange Kante kann dann folgendermaßen beschrieben werden:
Startknoten \sphinxhyphen{} d\_5\_10\_1 \sphinxhyphen{} d\_5\_10\_2 \sphinxhyphen{} … \sphinxhyphen{} d\_5\_10\_(span\sphinxhyphen{}1) \sphinxhyphen{} Endknoten. Außerdem werden die inter axis Knoten (Span = 0) ermittelt und gefiltert, um später die Achsenreihenfolge sauber zu ermitteln und die Pipelineschritte zu Initialisieren. Gefiltert bedeutet hier:
1. nur intra axis die nicht mixed sind (vgl. nested: is\_mixed, werden in 3b behandelt)
2. mixed (inter + intra, werden nur in 3a nicht aber in 3b berücksichtigt)
Reine intra\sphinxhyphen{}axis Knoten werden aus den regulären node\_groups entfernt und ihre Kanten vorübergehend aus dem Graphen gelöscht. Am Ende der Pipeline werden alle Informationen wieder konsistent zusammengesetzt.
\begin{quote}\begin{description}
\sphinxlineitem{Parameter}\begin{itemize}
\item {} 
\sphinxAtStartPar
\sphinxstyleliteralstrong{\sphinxupquote{layout}} (\DUrole{sphinx_autodoc_typehints-type}{\sphinxcode{\sphinxupquote{HivePlotLayout}}}) \textendash{} Das HivePlotLayout\sphinxhyphen{}Objekt, das die Informationen über den Graphen, die Achsen und die Knoten enthält.

\item {} 
\sphinxAtStartPar
\sphinxstyleliteralstrong{\sphinxupquote{node\_position\_map}} (\DUrole{sphinx_autodoc_typehints-type}{\sphinxcode{\sphinxupquote{dict}}{[}\sphinxcode{\sphinxupquote{int}} | \sphinxcode{\sphinxupquote{str}}, \sphinxcode{\sphinxupquote{int}}{]}}) \textendash{} Knoten\sphinxhyphen{}ID: Achsenindex in phi.

\item {} 
\sphinxAtStartPar
\sphinxstyleliteralstrong{\sphinxupquote{node\_axis\_map}} (\DUrole{sphinx_autodoc_typehints-type}{\sphinxcode{\sphinxupquote{dict}}{[}\sphinxcode{\sphinxupquote{int}} | \sphinxcode{\sphinxupquote{str}}, \sphinxcode{\sphinxupquote{int}}{]}}) \textendash{} Knoten\sphinxhyphen{}ID: Achse

\item {} 
\sphinxAtStartPar
\sphinxstyleliteralstrong{\sphinxupquote{(}}\sphinxstyleliteralstrong{\sphinxupquote{neighborhood\_map}} (\sphinxstyleliteralemphasis{\sphinxupquote{neighborhood\_map}}) \textendash{} dict{[}int | str, list{[}int | str{]}{]}): Knoten\sphinxhyphen{}ID: Nachbarknoten (:= proper = Span = 1)

\end{itemize}

\sphinxlineitem{Rückgabetyp}
\sphinxAtStartPar
\DUrole{sphinx_autodoc_typehints-type}{\sphinxcode{\sphinxupquote{None}}}

\sphinxlineitem{Rückgabe}
\sphinxAtStartPar
None. Die Funktion arbeitet in\sphinxhyphen{}place auf dem übergebenen Layout

\end{description}\end{quote}

\end{fulllineitems}



\section{src.dblp\_parser module}
\label{\detokenize{api:module-src.dblp_parser}}\label{\detokenize{api:src-dblp-parser-module}}\index{module@\spxentry{module}!src.dblp\_parser@\spxentry{src.dblp\_parser}}\index{src.dblp\_parser@\spxentry{src.dblp\_parser}!module@\spxentry{module}}\index{LocalDTDResolver (Klasse in src.dblp\_parser)@\spxentry{LocalDTDResolver}\spxextra{Klasse in src.dblp\_parser}}

\begin{fulllineitems}
\phantomsection\label{\detokenize{api:src.dblp_parser.LocalDTDResolver}}
\pysigstartsignatures
\pysiglinewithargsret
{\sphinxbfcode{\sphinxupquote{\DUrole{k}{class}\DUrole{w}{ }}}\sphinxcode{\sphinxupquote{src.dblp\_parser.}}\sphinxbfcode{\sphinxupquote{LocalDTDResolver}}}
{\sphinxparam{\DUrole{n}{dtd\_path}}}
{}
\pysigstopsignatures
\sphinxAtStartPar
Bases: \sphinxcode{\sphinxupquote{Resolver}}

\sphinxAtStartPar
Da es beim Parse trotz schneller Hardware zu Performanzproblemen kam mussten Einstellungen feingranular festgelegt werden. Unter anderem in dieser Klasse gebündelt die dtd\sphinxhyphen{}Datei zur Übersetzung von Umlauten und Sonderzeichen. Sie wird, falls vorhanden, von lokal bevorzugt.
\index{resolve() (Methode von src.dblp\_parser.LocalDTDResolver)@\spxentry{resolve()}\spxextra{Methode von src.dblp\_parser.LocalDTDResolver}}

\begin{fulllineitems}
\phantomsection\label{\detokenize{api:src.dblp_parser.LocalDTDResolver.resolve}}
\pysigstartsignatures
\pysiglinewithargsret
{\sphinxbfcode{\sphinxupquote{resolve}}}
{\sphinxparam{\DUrole{n}{self}}\sphinxparamcomma \sphinxparam{\DUrole{n}{system\_url}}\sphinxparamcomma \sphinxparam{\DUrole{n}{public\_id}}\sphinxparamcomma \sphinxparam{\DUrole{n}{context}}}
{}
\pysigstopsignatures
\sphinxAtStartPar
Override this method to resolve an external source by
\sphinxcode{\sphinxupquote{system\_url}} and \sphinxcode{\sphinxupquote{public\_id}}.  The third argument is an
opaque context object.

\sphinxAtStartPar
Return the result of one of the \sphinxcode{\sphinxupquote{resolve\_*()}} methods.

\end{fulllineitems}


\end{fulllineitems}

\index{build\_node\_identity\_maps() (im Modul src.dblp\_parser)@\spxentry{build\_node\_identity\_maps()}\spxextra{im Modul src.dblp\_parser}}

\begin{fulllineitems}
\phantomsection\label{\detokenize{api:src.dblp_parser.build_node_identity_maps}}
\pysigstartsignatures
\pysiglinewithargsret
{\sphinxcode{\sphinxupquote{src.dblp\_parser.}}\sphinxbfcode{\sphinxupquote{build\_node\_identity\_maps}}}
{\sphinxparam{\DUrole{n}{nodes}}}
{}
\pysigstopsignatures
\sphinxAtStartPar
Funktion dient der Erstellung der Maps für die gleichnamigen Felder im Hiveplotlayout. Es werden KnotenIDs als Integer auf die Namen der Forschenden als Strings und vice versa gemappt. Dadurch wird es einerseits ermöglicht den Eingabegraphen in ein pipelinefreundliches Format zu übersetzen und
gleichzeitig aus diesem am Ende die Label für den Renderer abzuleiten.
\begin{quote}\begin{description}
\sphinxlineitem{Rückgabetyp}
\sphinxAtStartPar
\DUrole{sphinx_autodoc_typehints-type}{\sphinxcode{\sphinxupquote{tuple}}{[}\sphinxcode{\sphinxupquote{dict}}{[}\sphinxcode{\sphinxupquote{int}}, \sphinxcode{\sphinxupquote{str}}{]}, \sphinxcode{\sphinxupquote{dict}}{[}\sphinxcode{\sphinxupquote{str}}, \sphinxcode{\sphinxupquote{int}}{]}{]}}

\end{description}\end{quote}

\end{fulllineitems}

\index{parse\_gd\_coauthor\_networks() (im Modul src.dblp\_parser)@\spxentry{parse\_gd\_coauthor\_networks()}\spxextra{im Modul src.dblp\_parser}}

\begin{fulllineitems}
\phantomsection\label{\detokenize{api:src.dblp_parser.parse_gd_coauthor_networks}}
\pysigstartsignatures
\pysiglinewithargsret
{\sphinxcode{\sphinxupquote{src.dblp\_parser.}}\sphinxbfcode{\sphinxupquote{parse\_gd\_coauthor\_networks}}}
{\sphinxparam{\DUrole{n}{xml\_path}}\sphinxparamcomma \sphinxparam{\DUrole{n}{cache\_path}\DUrole{o}{=}\DUrole{default_value}{None}}}
{}
\pysigstopsignatures
\sphinxAtStartPar
Parst dblp und gibt Co\sphinxhyphen{}Autoren\sphinxhyphen{}Graphen pro Jahr zurück. Ausgabe erfolgt als pickle, um die Daten wiederverwendbar zu halten, ohne den Parse wieder durchzuführen. Lädt aus Cache falls vorhanden.
\begin{quote}\begin{description}
\sphinxlineitem{Rückgabetyp}
\sphinxAtStartPar
\DUrole{sphinx_autodoc_typehints-type}{\sphinxcode{\sphinxupquote{dict}}{[}\sphinxcode{\sphinxupquote{int}}, \sphinxcode{\sphinxupquote{Graph}}{]}}

\end{description}\end{quote}

\end{fulllineitems}

\index{scan\_booktitles() (im Modul src.dblp\_parser)@\spxentry{scan\_booktitles()}\spxextra{im Modul src.dblp\_parser}}

\begin{fulllineitems}
\phantomsection\label{\detokenize{api:src.dblp_parser.scan_booktitles}}
\pysigstartsignatures
\pysiglinewithargsret
{\sphinxcode{\sphinxupquote{src.dblp\_parser.}}\sphinxbfcode{\sphinxupquote{scan\_booktitles}}}
{\sphinxparam{\DUrole{n}{xml\_path}}\sphinxparamcomma \sphinxparam{\DUrole{n}{dtd\_path}}}
{}
\pysigstopsignatures
\sphinxAtStartPar
Vorverarbeitung der dblp\sphinxhyphen{}xml Datei um nach booktitles zu filtern, da die xml\sphinxhyphen{}Datei neben Konferenzbänden auch andere wissenschaftliche Erzeugnisse wie Dissertationen oder Buchpublikationen enthält, die nichts mit der GD\sphinxhyphen{}Konferenz zu tun haben.
\begin{quote}\begin{description}
\sphinxlineitem{Rückgabetyp}
\sphinxAtStartPar
\DUrole{sphinx_autodoc_typehints-type}{\sphinxcode{\sphinxupquote{set}}{[}\sphinxcode{\sphinxupquote{str}}{]}}

\end{description}\end{quote}

\end{fulllineitems}



\section{src.experiments module}
\label{\detokenize{api:module-src.experiments}}\label{\detokenize{api:src-experiments-module}}\index{module@\spxentry{module}!src.experiments@\spxentry{src.experiments}}\index{src.experiments@\spxentry{src.experiments}!module@\spxentry{module}}\index{DataCollector (Klasse in src.experiments)@\spxentry{DataCollector}\spxextra{Klasse in src.experiments}}

\begin{fulllineitems}
\phantomsection\label{\detokenize{api:src.experiments.DataCollector}}
\pysigstartsignatures
\pysiglinewithargsret
{\sphinxbfcode{\sphinxupquote{\DUrole{k}{class}\DUrole{w}{ }}}\sphinxcode{\sphinxupquote{src.experiments.}}\sphinxbfcode{\sphinxupquote{DataCollector}}}
{\sphinxparam{\DUrole{n}{path}\DUrole{o}{=}\DUrole{default_value}{None}}}
{}
\pysigstopsignatures
\sphinxAtStartPar
Bases: \sphinxcode{\sphinxupquote{object}}

\sphinxAtStartPar
Zentrale Datenstruktur für den Datensatz der Experimente.
\index{data\_keys() (Methode von src.experiments.DataCollector)@\spxentry{data\_keys()}\spxextra{Methode von src.experiments.DataCollector}}

\begin{fulllineitems}
\phantomsection\label{\detokenize{api:src.experiments.DataCollector.data_keys}}
\pysigstartsignatures
\pysiglinewithargsret
{\sphinxbfcode{\sphinxupquote{data\_keys}}}
{\sphinxparam{\DUrole{n}{data\_set}\DUrole{o}{=}\DUrole{default_value}{None}}}
{}
\pysigstopsignatures\begin{quote}\begin{description}
\sphinxlineitem{Rückgabetyp}
\sphinxAtStartPar
\DUrole{sphinx_autodoc_typehints-type}{\sphinxcode{\sphinxupquote{list}}{[}\sphinxcode{\sphinxupquote{str}}{]}}

\end{description}\end{quote}

\end{fulllineitems}

\index{data\_sets (Attribut von src.experiments.DataCollector)@\spxentry{data\_sets}\spxextra{Attribut von src.experiments.DataCollector}}

\begin{fulllineitems}
\phantomsection\label{\detokenize{api:src.experiments.DataCollector.data_sets}}
\pysigstartsignatures
\pysigline
{\sphinxbfcode{\sphinxupquote{data\_sets}}\sphinxbfcode{\sphinxupquote{\DUrole{p}{:}\DUrole{w}{ }dict\DUrole{p}{{[}}str\DUrole{p}{,}\DUrole{w}{ }dict\DUrole{p}{{[}}str\DUrole{p}{,}\DUrole{w}{ }list\DUrole{p}{{]}}\DUrole{p}{{]}}}}}
\pysigstopsignatures
\end{fulllineitems}

\index{delete() (Methode von src.experiments.DataCollector)@\spxentry{delete()}\spxextra{Methode von src.experiments.DataCollector}}

\begin{fulllineitems}
\phantomsection\label{\detokenize{api:src.experiments.DataCollector.delete}}
\pysigstartsignatures
\pysiglinewithargsret
{\sphinxbfcode{\sphinxupquote{delete}}}
{\sphinxparam{\DUrole{n}{data\_set}}\sphinxparamcomma \sphinxparam{\DUrole{n}{x\_val}\DUrole{o}{=}\DUrole{default_value}{None}}}
{}
\pysigstopsignatures\begin{quote}\begin{description}
\sphinxlineitem{Rückgabetyp}
\sphinxAtStartPar
\DUrole{sphinx_autodoc_typehints-type}{\sphinxcode{\sphinxupquote{None}}}

\end{description}\end{quote}

\end{fulllineitems}

\index{get\_data\_set() (Methode von src.experiments.DataCollector)@\spxentry{get\_data\_set()}\spxextra{Methode von src.experiments.DataCollector}}

\begin{fulllineitems}
\phantomsection\label{\detokenize{api:src.experiments.DataCollector.get_data_set}}
\pysigstartsignatures
\pysiglinewithargsret
{\sphinxbfcode{\sphinxupquote{get\_data\_set}}}
{\sphinxparam{\DUrole{n}{data\_set}}}
{}
\pysigstopsignatures
\end{fulllineitems}

\index{merge() (Methode von src.experiments.DataCollector)@\spxentry{merge()}\spxextra{Methode von src.experiments.DataCollector}}

\begin{fulllineitems}
\phantomsection\label{\detokenize{api:src.experiments.DataCollector.merge}}
\pysigstartsignatures
\pysiglinewithargsret
{\sphinxbfcode{\sphinxupquote{merge}}}
{}
{}
\pysigstopsignatures
\end{fulllineitems}

\index{missing() (Methode von src.experiments.DataCollector)@\spxentry{missing()}\spxextra{Methode von src.experiments.DataCollector}}

\begin{fulllineitems}
\phantomsection\label{\detokenize{api:src.experiments.DataCollector.missing}}
\pysigstartsignatures
\pysiglinewithargsret
{\sphinxbfcode{\sphinxupquote{missing}}}
{}
{}
\pysigstopsignatures
\sphinxAtStartPar
Überprüft die Datensätze auf fehlende Einträge. Falls welche gefunden werden gibt es eine Liste aus Tupeln zurück,
wobei das Tupel (Datensatz, x, fehlende Datenpunktbezeichner) ist. Falls x schon fehlt wird dies kenntlich gemacht.
\begin{quote}\begin{description}
\sphinxlineitem{Rückgabe}
\sphinxAtStartPar
Fehlende Einträge

\sphinxlineitem{Rückgabetyp}
\sphinxAtStartPar
\DUrole{sphinx_autodoc_typehints-type}{\sphinxcode{\sphinxupquote{list}}{[}\sphinxcode{\sphinxupquote{tuple}}{[}\sphinxcode{\sphinxupquote{str}}, \sphinxcode{\sphinxupquote{str}} | \sphinxcode{\sphinxupquote{int}} | \sphinxcode{\sphinxupquote{float}}, \sphinxcode{\sphinxupquote{list}}{[}\sphinxcode{\sphinxupquote{str}}{]}{]}{]}}

\end{description}\end{quote}

\end{fulllineitems}

\index{save() (Methode von src.experiments.DataCollector)@\spxentry{save()}\spxextra{Methode von src.experiments.DataCollector}}

\begin{fulllineitems}
\phantomsection\label{\detokenize{api:src.experiments.DataCollector.save}}
\pysigstartsignatures
\pysiglinewithargsret
{\sphinxbfcode{\sphinxupquote{save}}}
{}
{}
\pysigstopsignatures\begin{quote}\begin{description}
\sphinxlineitem{Rückgabetyp}
\sphinxAtStartPar
\DUrole{sphinx_autodoc_typehints-type}{\sphinxcode{\sphinxupquote{None}}}

\end{description}\end{quote}

\end{fulllineitems}

\index{update() (Methode von src.experiments.DataCollector)@\spxentry{update()}\spxextra{Methode von src.experiments.DataCollector}}

\begin{fulllineitems}
\phantomsection\label{\detokenize{api:src.experiments.DataCollector.update}}
\pysigstartsignatures
\pysiglinewithargsret
{\sphinxbfcode{\sphinxupquote{update}}}
{\sphinxparam{\DUrole{n}{data\_set}}\sphinxparamcomma \sphinxparam{\DUrole{n}{x\_val}\DUrole{o}{=}\DUrole{default_value}{None}}\sphinxparamcomma \sphinxparam{\DUrole{o}{**}\DUrole{n}{kwargs}}}
{}
\pysigstopsignatures\begin{quote}\begin{description}
\sphinxlineitem{Rückgabetyp}
\sphinxAtStartPar
\DUrole{sphinx_autodoc_typehints-type}{\sphinxcode{\sphinxupquote{None}}}

\end{description}\end{quote}

\end{fulllineitems}

\index{validate() (Methode von src.experiments.DataCollector)@\spxentry{validate()}\spxextra{Methode von src.experiments.DataCollector}}

\begin{fulllineitems}
\phantomsection\label{\detokenize{api:src.experiments.DataCollector.validate}}
\pysigstartsignatures
\pysiglinewithargsret
{\sphinxbfcode{\sphinxupquote{validate}}}
{\sphinxparam{\DUrole{n}{data\_set}}}
{}
\pysigstopsignatures
\sphinxAtStartPar
Ermittelt ob die Einträge des Datensatzes gleich viele Datenpunkte enthalten.
\begin{quote}\begin{description}
\sphinxlineitem{Parameter}
\sphinxAtStartPar
\sphinxstyleliteralstrong{\sphinxupquote{data\_set}} (\DUrole{sphinx_autodoc_typehints-type}{\sphinxcode{\sphinxupquote{str}}}) \textendash{} betrachteter Datensatz

\sphinxlineitem{Rückgabe}
\sphinxAtStartPar
True = wenn Einträge gleiche Länge, sonst False. falls Datensatz fehlt Debug\sphinxhyphen{}Ausgabe

\sphinxlineitem{Rückgabetyp}
\sphinxAtStartPar
\DUrole{sphinx_autodoc_typehints-type}{\sphinxcode{\sphinxupquote{bool}} | \sphinxcode{\sphinxupquote{None}}}

\end{description}\end{quote}

\end{fulllineitems}


\end{fulllineitems}



\section{src.hiveplot module}
\label{\detokenize{api:module-src.hiveplot}}\label{\detokenize{api:src-hiveplot-module}}\index{module@\spxentry{module}!src.hiveplot@\spxentry{src.hiveplot}}\index{src.hiveplot@\spxentry{src.hiveplot}!module@\spxentry{module}}\index{HivePlotLayout (Klasse in src.hiveplot)@\spxentry{HivePlotLayout}\spxextra{Klasse in src.hiveplot}}

\begin{fulllineitems}
\phantomsection\label{\detokenize{api:src.hiveplot.HivePlotLayout}}
\pysigstartsignatures
\pysiglinewithargsret
{\sphinxbfcode{\sphinxupquote{\DUrole{k}{class}\DUrole{w}{ }}}\sphinxcode{\sphinxupquote{src.hiveplot.}}\sphinxbfcode{\sphinxupquote{HivePlotLayout}}}
{\sphinxparam{\DUrole{n}{graph}}\sphinxparamcomma \sphinxparam{\DUrole{n}{num\_axes}}\sphinxparamcomma \sphinxparam{\DUrole{n}{axis\_order}}\sphinxparamcomma \sphinxparam{\DUrole{n}{node\_groups}}\sphinxparamcomma \sphinxparam{\DUrole{n}{node\_groups\_expanded}\DUrole{o}{=}\DUrole{default_value}{\textless{}factory\textgreater{}}}\sphinxparamcomma \sphinxparam{\DUrole{n}{edges\_expanded}\DUrole{o}{=}\DUrole{default_value}{\textless{}factory\textgreater{}}}\sphinxparamcomma \sphinxparam{\DUrole{n}{node\_groups\_dummies}\DUrole{o}{=}\DUrole{default_value}{\textless{}factory\textgreater{}}}\sphinxparamcomma \sphinxparam{\DUrole{n}{dummy\_edge\_segments}\DUrole{o}{=}\DUrole{default_value}{\textless{}factory\textgreater{}}}\sphinxparamcomma \sphinxparam{\DUrole{n}{long\_edges}\DUrole{o}{=}\DUrole{default_value}{\textless{}factory\textgreater{}}}\sphinxparamcomma \sphinxparam{\DUrole{n}{intra\_axis\_nodes}\DUrole{o}{=}\DUrole{default_value}{\textless{}factory\textgreater{}}}\sphinxparamcomma \sphinxparam{\DUrole{n}{intra\_axis\_edges}\DUrole{o}{=}\DUrole{default_value}{\textless{}factory\textgreater{}}}\sphinxparamcomma \sphinxparam{\DUrole{n}{id\_to\_name}\DUrole{o}{=}\DUrole{default_value}{\textless{}factory\textgreater{}}}\sphinxparamcomma \sphinxparam{\DUrole{n}{name\_to\_id}\DUrole{o}{=}\DUrole{default_value}{\textless{}factory\textgreater{}}}\sphinxparamcomma \sphinxparam{\DUrole{n}{crossings\_expanded}\DUrole{o}{=}\DUrole{default_value}{None}}\sphinxparamcomma \sphinxparam{\DUrole{n}{mixed\_nodes\_by\_axis}\DUrole{o}{=}\DUrole{default_value}{\textless{}factory\textgreater{}}}\sphinxparamcomma \sphinxparam{\DUrole{n}{strict\_intra\_nodes\_by\_axis}\DUrole{o}{=}\DUrole{default_value}{\textless{}factory\textgreater{}}}\sphinxparamcomma \sphinxparam{\DUrole{n}{fixed\_inter\_axis\_delta}\DUrole{o}{=}\DUrole{default_value}{None}}\sphinxparamcomma \sphinxparam{\DUrole{n}{fixed\_positions\_by\_axis}\DUrole{o}{=}\DUrole{default_value}{None}}}
{}
\pysigstopsignatures
\sphinxAtStartPar
Bases: \sphinxcode{\sphinxupquote{object}}

\sphinxAtStartPar
Zentrale Datenstruktur für ein berechnetes Hive\sphinxhyphen{}Plot\sphinxhyphen{}Layout.
\index{graph (Attribut von src.hiveplot.HivePlotLayout)@\spxentry{graph}\spxextra{Attribut von src.hiveplot.HivePlotLayout}}

\begin{fulllineitems}
\phantomsection\label{\detokenize{api:src.hiveplot.HivePlotLayout.graph}}
\pysigstartsignatures
\pysigline
{\sphinxbfcode{\sphinxupquote{graph}}}
\pysigstopsignatures
\sphinxAtStartPar
ein Simple Graph den man nach der Verarbeitung der geparsten DBLP Daten erhält: V = Autoren, E = Coautorenschaft zu
\begin{quote}\begin{description}
\sphinxlineitem{Type}
\sphinxAtStartPar
nx.Graph

\end{description}\end{quote}

\end{fulllineitems}

\index{num\_axes (Attribut von src.hiveplot.HivePlotLayout)@\spxentry{num\_axes}\spxextra{Attribut von src.hiveplot.HivePlotLayout}}

\begin{fulllineitems}
\phantomsection\label{\detokenize{api:src.hiveplot.HivePlotLayout.num_axes}}
\pysigstartsignatures
\pysigline
{\sphinxbfcode{\sphinxupquote{num\_axes}}}
\pysigstopsignatures
\sphinxAtStartPar
k = Anzahl Achsen
\begin{quote}\begin{description}
\sphinxlineitem{Type}
\sphinxAtStartPar
int

\end{description}\end{quote}

\end{fulllineitems}

\index{axis\_order (Attribut von src.hiveplot.HivePlotLayout)@\spxentry{axis\_order}\spxextra{Attribut von src.hiveplot.HivePlotLayout}}

\begin{fulllineitems}
\phantomsection\label{\detokenize{api:src.hiveplot.HivePlotLayout.axis_order}}
\pysigstartsignatures
\pysigline
{\sphinxbfcode{\sphinxupquote{axis\_order}}}
\pysigstopsignatures
\sphinxAtStartPar
phi
\begin{quote}\begin{description}
\sphinxlineitem{Type}
\sphinxAtStartPar
list{[}int{]}

\end{description}\end{quote}

\end{fulllineitems}

\index{node\_groups (Attribut von src.hiveplot.HivePlotLayout)@\spxentry{node\_groups}\spxextra{Attribut von src.hiveplot.HivePlotLayout}}

\begin{fulllineitems}
\phantomsection\label{\detokenize{api:src.hiveplot.HivePlotLayout.node_groups}}
\pysigstartsignatures
\pysigline
{\sphinxbfcode{\sphinxupquote{node\_groups}}}
\pysigstopsignatures
\sphinxAtStartPar
alpha: siehe node\_groups Funktion in ordering.py
\begin{quote}\begin{description}
\sphinxlineitem{Type}
\sphinxAtStartPar
dict{[}int, list{[}int{]}{]}

\end{description}\end{quote}

\end{fulllineitems}

\index{node\_groups\_dummies (Attribut von src.hiveplot.HivePlotLayout)@\spxentry{node\_groups\_dummies}\spxextra{Attribut von src.hiveplot.HivePlotLayout}}

\begin{fulllineitems}
\phantomsection\label{\detokenize{api:src.hiveplot.HivePlotLayout.node_groups_dummies}}
\pysigstartsignatures
\pysigline
{\sphinxbfcode{\sphinxupquote{node\_groups\_dummies}}}
\pysigstopsignatures
\sphinxAtStartPar
Achsen mit Dummyknoten
\begin{quote}\begin{description}
\sphinxlineitem{Type}
\sphinxAtStartPar
dict{[}int, list{[}int{]}{]}

\end{description}\end{quote}

\end{fulllineitems}



\begin{fulllineitems}

\pysigstartsignatures
\pysigline
{\sphinxbfcode{\sphinxupquote{\#~dummy\_edge\_segments}}}
\pysigstopsignatures
\sphinxAtStartPar
key: Kantentupel (u,v) value: Liste von Kantensegmenten die u und v über Dummyknoten verbinden
\begin{quote}\begin{description}
\sphinxlineitem{Type}
\sphinxAtStartPar
dict{[}tuple{[}int,int{]}, list{[}tuple{[}int, int{]}{]}{]}

\end{description}\end{quote}

\end{fulllineitems}

\index{node\_order (Attribut von src.hiveplot.HivePlotLayout)@\spxentry{node\_order}\spxextra{Attribut von src.hiveplot.HivePlotLayout}}

\begin{fulllineitems}
\phantomsection\label{\detokenize{api:src.hiveplot.HivePlotLayout.node_order}}
\pysigstartsignatures
\pysigline
{\sphinxbfcode{\sphinxupquote{node\_order}}}
\pysigstopsignatures
\sphinxAtStartPar
pi\_i pro Achse gebündelt
\begin{quote}\begin{description}
\sphinxlineitem{Type}
\sphinxAtStartPar
dict{[}int, list{[}int{]}{]}

\end{description}\end{quote}

\end{fulllineitems}

\index{crossings (Attribut von src.hiveplot.HivePlotLayout)@\spxentry{crossings}\spxextra{Attribut von src.hiveplot.HivePlotLayout}}

\begin{fulllineitems}
\phantomsection\label{\detokenize{api:src.hiveplot.HivePlotLayout.crossings}}
\pysigstartsignatures
\pysigline
{\sphinxbfcode{\sphinxupquote{crossings}}}
\pysigstopsignatures
\sphinxAtStartPar
Kreuzungszahl Standardmodell
\begin{quote}\begin{description}
\sphinxlineitem{Type}
\sphinxAtStartPar
Optional{[}int{]}

\end{description}\end{quote}

\end{fulllineitems}

\index{crossings\_extended (Attribut von src.hiveplot.HivePlotLayout)@\spxentry{crossings\_extended}\spxextra{Attribut von src.hiveplot.HivePlotLayout}}

\begin{fulllineitems}
\phantomsection\label{\detokenize{api:src.hiveplot.HivePlotLayout.crossings_extended}}
\pysigstartsignatures
\pysigline
{\sphinxbfcode{\sphinxupquote{crossings\_extended}}}
\pysigstopsignatures
\sphinxAtStartPar
Kreuzungszahl erweitertes Modell
\begin{quote}\begin{description}
\sphinxlineitem{Type}
\sphinxAtStartPar
Optional{[}int{]}

\end{description}\end{quote}

\end{fulllineitems}

\index{axis\_order (Attribut von src.hiveplot.HivePlotLayout)@\spxentry{axis\_order}\spxextra{Attribut von src.hiveplot.HivePlotLayout}}

\begin{fulllineitems}
\phantomsection\label{\detokenize{api:id0}}
\pysigstartsignatures
\pysigline
{\sphinxbfcode{\sphinxupquote{axis\_order}}\sphinxbfcode{\sphinxupquote{\DUrole{p}{:}\DUrole{w}{ }list\DUrole{p}{{[}}int\DUrole{p}{{]}}}}}
\pysigstopsignatures
\end{fulllineitems}

\index{classify\_nodes\_for\_3b() (Methode von src.hiveplot.HivePlotLayout)@\spxentry{classify\_nodes\_for\_3b()}\spxextra{Methode von src.hiveplot.HivePlotLayout}}

\begin{fulllineitems}
\phantomsection\label{\detokenize{api:src.hiveplot.HivePlotLayout.classify_nodes_for_3b}}
\pysigstartsignatures
\pysiglinewithargsret
{\sphinxbfcode{\sphinxupquote{classify\_nodes\_for\_3b}}}
{}
{}
\pysigstopsignatures
\sphinxAtStartPar
Klassifiziert die Knoten für Pipelineschritt 3b.
\begin{quote}\begin{description}
\sphinxlineitem{Rückgabetyp}
\sphinxAtStartPar
\DUrole{sphinx_autodoc_typehints-type}{\sphinxcode{\sphinxupquote{None}}}

\end{description}\end{quote}

\end{fulllineitems}

\index{copy() (Methode von src.hiveplot.HivePlotLayout)@\spxentry{copy()}\spxextra{Methode von src.hiveplot.HivePlotLayout}}

\begin{fulllineitems}
\phantomsection\label{\detokenize{api:src.hiveplot.HivePlotLayout.copy}}
\pysigstartsignatures
\pysiglinewithargsret
{\sphinxbfcode{\sphinxupquote{copy}}}
{}
{}
\pysigstopsignatures
\sphinxAtStartPar
Gibt eine Deepcopy des Hiveplotlayouts zurück.
\begin{quote}\begin{description}
\sphinxlineitem{Rückgabetyp}
\sphinxAtStartPar
\DUrole{sphinx_autodoc_typehints-type}{{\hyperref[\detokenize{api:src.hiveplot.HivePlotLayout.copy}]{\sphinxcrossref{\sphinxcode{\sphinxupquote{copy}}}}}}

\end{description}\end{quote}

\end{fulllineitems}

\index{count\_crossings() (Methode von src.hiveplot.HivePlotLayout)@\spxentry{count\_crossings()}\spxextra{Methode von src.hiveplot.HivePlotLayout}}

\begin{fulllineitems}
\phantomsection\label{\detokenize{api:src.hiveplot.HivePlotLayout.count_crossings}}
\pysigstartsignatures
\pysiglinewithargsret
{\sphinxbfcode{\sphinxupquote{count\_crossings}}}
{\sphinxparam{\DUrole{n}{layout\_expanded}\DUrole{o}{=}\DUrole{default_value}{False}}}
{}
\pysigstopsignatures
\sphinxAtStartPar
Funktion zählt die induzierten Kanten des fertigen Hiveplotlayout.
\begin{quote}\begin{description}
\sphinxlineitem{Rückgabetyp}
\sphinxAtStartPar
\DUrole{sphinx_autodoc_typehints-type}{\sphinxcode{\sphinxupquote{int}}}

\end{description}\end{quote}

\end{fulllineitems}

\index{crossings\_expanded (Attribut von src.hiveplot.HivePlotLayout)@\spxentry{crossings\_expanded}\spxextra{Attribut von src.hiveplot.HivePlotLayout}}

\begin{fulllineitems}
\phantomsection\label{\detokenize{api:src.hiveplot.HivePlotLayout.crossings_expanded}}
\pysigstartsignatures
\pysigline
{\sphinxbfcode{\sphinxupquote{crossings\_expanded}}\sphinxbfcode{\sphinxupquote{\DUrole{p}{:}\DUrole{w}{ }int\DUrole{w}{ }\DUrole{p}{|}\DUrole{w}{ }None}}\sphinxbfcode{\sphinxupquote{\DUrole{w}{ }\DUrole{p}{=}\DUrole{w}{ }None}}}
\pysigstopsignatures
\end{fulllineitems}

\index{dummy\_edge\_segments (Attribut von src.hiveplot.HivePlotLayout)@\spxentry{dummy\_edge\_segments}\spxextra{Attribut von src.hiveplot.HivePlotLayout}}

\begin{fulllineitems}
\phantomsection\label{\detokenize{api:src.hiveplot.HivePlotLayout.dummy_edge_segments}}
\pysigstartsignatures
\pysigline
{\sphinxbfcode{\sphinxupquote{dummy\_edge\_segments}}\sphinxbfcode{\sphinxupquote{\DUrole{p}{:}\DUrole{w}{ }list\DUrole{p}{{[}}tuple\DUrole{p}{{[}}int\DUrole{w}{ }\DUrole{p}{|}\DUrole{w}{ }str\DUrole{p}{,}\DUrole{w}{ }int\DUrole{w}{ }\DUrole{p}{|}\DUrole{w}{ }str\DUrole{p}{{]}}\DUrole{p}{{]}}}}}
\pysigstopsignatures
\end{fulllineitems}

\index{edges() (Methode von src.hiveplot.HivePlotLayout)@\spxentry{edges()}\spxextra{Methode von src.hiveplot.HivePlotLayout}}

\begin{fulllineitems}
\phantomsection\label{\detokenize{api:src.hiveplot.HivePlotLayout.edges}}
\pysigstartsignatures
\pysiglinewithargsret
{\sphinxbfcode{\sphinxupquote{edges}}}
{}
{}
\pysigstopsignatures
\sphinxAtStartPar
Gibt die Kantenliste des Graphen zurück. Notwendig, weil mit dem Networkx Edgeview nicht direkt gearbeitet werden kann.
\begin{quote}\begin{description}
\sphinxlineitem{Rückgabe}
\sphinxAtStartPar
Liste der Kanten als Tupel (u, v)

\sphinxlineitem{Rückgabetyp}
\sphinxAtStartPar
\DUrole{sphinx_autodoc_typehints-type}{\sphinxcode{\sphinxupquote{list}}{[}\sphinxcode{\sphinxupquote{tuple}}{]}}

\end{description}\end{quote}

\end{fulllineitems}

\index{edges\_expanded (Attribut von src.hiveplot.HivePlotLayout)@\spxentry{edges\_expanded}\spxextra{Attribut von src.hiveplot.HivePlotLayout}}

\begin{fulllineitems}
\phantomsection\label{\detokenize{api:src.hiveplot.HivePlotLayout.edges_expanded}}
\pysigstartsignatures
\pysigline
{\sphinxbfcode{\sphinxupquote{edges\_expanded}}\sphinxbfcode{\sphinxupquote{\DUrole{p}{:}\DUrole{w}{ }list\DUrole{p}{{[}}tuple\DUrole{p}{{[}}int\DUrole{w}{ }\DUrole{p}{|}\DUrole{w}{ }str\DUrole{p}{,}\DUrole{w}{ }int\DUrole{w}{ }\DUrole{p}{|}\DUrole{w}{ }str\DUrole{p}{{]}}\DUrole{p}{{]}}}}}
\pysigstopsignatures
\end{fulllineitems}

\index{fixed\_inter\_axis\_delta (Attribut von src.hiveplot.HivePlotLayout)@\spxentry{fixed\_inter\_axis\_delta}\spxextra{Attribut von src.hiveplot.HivePlotLayout}}

\begin{fulllineitems}
\phantomsection\label{\detokenize{api:src.hiveplot.HivePlotLayout.fixed_inter_axis_delta}}
\pysigstartsignatures
\pysigline
{\sphinxbfcode{\sphinxupquote{fixed\_inter\_axis\_delta}}\sphinxbfcode{\sphinxupquote{\DUrole{p}{:}\DUrole{w}{ }dict\DUrole{w}{ }\DUrole{p}{|}\DUrole{w}{ }None}}\sphinxbfcode{\sphinxupquote{\DUrole{w}{ }\DUrole{p}{=}\DUrole{w}{ }None}}}
\pysigstopsignatures
\end{fulllineitems}

\index{fixed\_positions\_by\_axis (Attribut von src.hiveplot.HivePlotLayout)@\spxentry{fixed\_positions\_by\_axis}\spxextra{Attribut von src.hiveplot.HivePlotLayout}}

\begin{fulllineitems}
\phantomsection\label{\detokenize{api:src.hiveplot.HivePlotLayout.fixed_positions_by_axis}}
\pysigstartsignatures
\pysigline
{\sphinxbfcode{\sphinxupquote{fixed\_positions\_by\_axis}}\sphinxbfcode{\sphinxupquote{\DUrole{p}{:}\DUrole{w}{ }dict\DUrole{w}{ }\DUrole{p}{|}\DUrole{w}{ }None}}\sphinxbfcode{\sphinxupquote{\DUrole{w}{ }\DUrole{p}{=}\DUrole{w}{ }None}}}
\pysigstopsignatures
\end{fulllineitems}

\index{freeze\_barycenter\_positions() (Methode von src.hiveplot.HivePlotLayout)@\spxentry{freeze\_barycenter\_positions()}\spxextra{Methode von src.hiveplot.HivePlotLayout}}

\begin{fulllineitems}
\phantomsection\label{\detokenize{api:src.hiveplot.HivePlotLayout.freeze_barycenter_positions}}
\pysigstartsignatures
\pysiglinewithargsret
{\sphinxbfcode{\sphinxupquote{freeze\_barycenter\_positions}}}
{\sphinxparam{\DUrole{n}{layout\_expanded}\DUrole{o}{=}\DUrole{default_value}{False}}}
{}
\pysigstopsignatures
\sphinxAtStartPar
Legt die fixierten Positionen für den Pipelineschritt 3b der Barycenterpipeline fest. Wird für alle mixed Knoten ermittelt.

\end{fulllineitems}

\index{freeze\_inter\_axis\_delta() (Methode von src.hiveplot.HivePlotLayout)@\spxentry{freeze\_inter\_axis\_delta()}\spxextra{Methode von src.hiveplot.HivePlotLayout}}

\begin{fulllineitems}
\phantomsection\label{\detokenize{api:src.hiveplot.HivePlotLayout.freeze_inter_axis_delta}}
\pysigstartsignatures
\pysiglinewithargsret
{\sphinxbfcode{\sphinxupquote{freeze\_inter\_axis\_delta}}}
{}
{}
\pysigstopsignatures
\sphinxAtStartPar
Legt die delta Map für den Pipelineschritt 3b der ILP\sphinxhyphen{}Pipeline fest. Wird für alle mixed Knoten ermittelt.
\begin{quote}\begin{description}
\sphinxlineitem{Rückgabetyp}
\sphinxAtStartPar
\DUrole{sphinx_autodoc_typehints-type}{\sphinxcode{\sphinxupquote{None}}}

\end{description}\end{quote}

\end{fulllineitems}

\index{fuse\_edges\_with\_edge\_dummies() (Methode von src.hiveplot.HivePlotLayout)@\spxentry{fuse\_edges\_with\_edge\_dummies()}\spxextra{Methode von src.hiveplot.HivePlotLayout}}

\begin{fulllineitems}
\phantomsection\label{\detokenize{api:src.hiveplot.HivePlotLayout.fuse_edges_with_edge_dummies}}
\pysigstartsignatures
\pysiglinewithargsret
{\sphinxbfcode{\sphinxupquote{fuse\_edges\_with\_edge\_dummies}}}
{}
{}
\pysigstopsignatures
\sphinxAtStartPar
Erzeugt eine Liste die alle kurzen Kanten und Dummykanten vereinigt zurückgibt.
\begin{quote}\begin{description}
\sphinxlineitem{Rückgabe}
\sphinxAtStartPar
Vereinigung aus Kanten und Dummykanten

\sphinxlineitem{Rückgabetyp}
\sphinxAtStartPar
\DUrole{sphinx_autodoc_typehints-type}{\sphinxcode{\sphinxupquote{list}}{[}\sphinxcode{\sphinxupquote{tuple}}{[}\sphinxcode{\sphinxupquote{int}} | \sphinxcode{\sphinxupquote{str}}, \sphinxcode{\sphinxupquote{int}} | \sphinxcode{\sphinxupquote{str}}{]}{]}}

\end{description}\end{quote}

\end{fulllineitems}

\index{fuse\_node\_groups\_with\_dummies() (Methode von src.hiveplot.HivePlotLayout)@\spxentry{fuse\_node\_groups\_with\_dummies()}\spxextra{Methode von src.hiveplot.HivePlotLayout}}

\begin{fulllineitems}
\phantomsection\label{\detokenize{api:src.hiveplot.HivePlotLayout.fuse_node_groups_with_dummies}}
\pysigstartsignatures
\pysiglinewithargsret
{\sphinxbfcode{\sphinxupquote{fuse\_node\_groups\_with\_dummies}}}
{\sphinxparam{\DUrole{n}{layout\_expanded}\DUrole{o}{=}\DUrole{default_value}{False}}}
{}
\pysigstopsignatures
\sphinxAtStartPar
Erzeugt ein dict mit Achsen als Schlüssel und der vereinigten Menge aus Knoten und Dummyknoten pro Achse als Wert. Zuerst die realen dann die virtuellen Knoten.
:param layout\_expanded: dient der Unterscheidung, ob in der Pipeline mit expandierten Achsen gerechnet wird oder nicht, Default = False (nicht expandierter Fall)
:type layout\_expanded: \DUrole{sphinx_autodoc_typehints-type}{\sphinxcode{\sphinxupquote{bool}}}
\begin{quote}\begin{description}
\sphinxlineitem{Rückgabe}
\sphinxAtStartPar
Vereinigung der Knoten und Dummyknoten pro Achse.

\sphinxlineitem{Rückgabetyp}
\sphinxAtStartPar
\DUrole{sphinx_autodoc_typehints-type}{\sphinxcode{\sphinxupquote{dict}}{[}\sphinxcode{\sphinxupquote{int}}, \sphinxcode{\sphinxupquote{list}}{[}\sphinxcode{\sphinxupquote{int}} | \sphinxcode{\sphinxupquote{str}}{]}{]}}

\end{description}\end{quote}

\end{fulllineitems}

\index{get\_proper\_neighborhood\_map() (Methode von src.hiveplot.HivePlotLayout)@\spxentry{get\_proper\_neighborhood\_map()}\spxextra{Methode von src.hiveplot.HivePlotLayout}}

\begin{fulllineitems}
\phantomsection\label{\detokenize{api:src.hiveplot.HivePlotLayout.get_proper_neighborhood_map}}
\pysigstartsignatures
\pysiglinewithargsret
{\sphinxbfcode{\sphinxupquote{get\_proper\_neighborhood\_map}}}
{\sphinxparam{\DUrole{n}{fused\_edge\_list}}\sphinxparamcomma \sphinxparam{\DUrole{n}{layout\_expanded}\DUrole{o}{=}\DUrole{default_value}{False}}}
{}
\pysigstopsignatures
\sphinxAtStartPar
Die Funktion ermittelt eine Liste, die jeden Knoten (sowohl real als auch virtuell) auf eine Liste seiner Nachbarn mappt.
\begin{quote}\begin{description}
\sphinxlineitem{Parameter}\begin{itemize}
\item {} 
\sphinxAtStartPar
\sphinxstyleliteralstrong{\sphinxupquote{fused\_edge\_list}} (\DUrole{sphinx_autodoc_typehints-type}{\sphinxcode{\sphinxupquote{list}}{[}\sphinxcode{\sphinxupquote{tuple}}{[}\sphinxcode{\sphinxupquote{int}} | \sphinxcode{\sphinxupquote{str}}, \sphinxcode{\sphinxupquote{int}} | \sphinxcode{\sphinxupquote{str}}{]}{]}}) \textendash{} list{[}tuple{[}int | str, int | str{]}{]}: Liste der realen und virtuellen Kantentupel

\item {} 
\sphinxAtStartPar
\sphinxstyleliteralstrong{\sphinxupquote{layout\_expanded}} (\DUrole{sphinx_autodoc_typehints-type}{\sphinxcode{\sphinxupquote{bool}}}) \textendash{} dient der Unterscheidung, ob in der Pipeline mit expandierten Achsen gerechnet wird oder nicht, Default = False (nicht expandierter Fall)

\end{itemize}

\sphinxlineitem{Rückgabe}
\sphinxAtStartPar
list{[}int | str{]}{]}: KnotenID: Nachbarliste

\sphinxlineitem{Rückgabetyp}
\sphinxAtStartPar
\DUrole{sphinx_autodoc_typehints-type}{\sphinxcode{\sphinxupquote{dict}}{[}slice(int | str, list{[}int | str{]}, None){]}}

\end{description}\end{quote}

\end{fulllineitems}

\index{graph (Attribut von src.hiveplot.HivePlotLayout)@\spxentry{graph}\spxextra{Attribut von src.hiveplot.HivePlotLayout}}

\begin{fulllineitems}
\phantomsection\label{\detokenize{api:id1}}
\pysigstartsignatures
\pysigline
{\sphinxbfcode{\sphinxupquote{graph}}\sphinxbfcode{\sphinxupquote{\DUrole{p}{:}\DUrole{w}{ }Graph}}}
\pysigstopsignatures
\end{fulllineitems}

\index{id\_to\_name (Attribut von src.hiveplot.HivePlotLayout)@\spxentry{id\_to\_name}\spxextra{Attribut von src.hiveplot.HivePlotLayout}}

\begin{fulllineitems}
\phantomsection\label{\detokenize{api:src.hiveplot.HivePlotLayout.id_to_name}}
\pysigstartsignatures
\pysigline
{\sphinxbfcode{\sphinxupquote{id\_to\_name}}\sphinxbfcode{\sphinxupquote{\DUrole{p}{:}\DUrole{w}{ }dict\DUrole{p}{{[}}int\DUrole{p}{,}\DUrole{w}{ }str\DUrole{p}{{]}}}}}
\pysigstopsignatures
\end{fulllineitems}

\index{intra\_axis\_edges (Attribut von src.hiveplot.HivePlotLayout)@\spxentry{intra\_axis\_edges}\spxextra{Attribut von src.hiveplot.HivePlotLayout}}

\begin{fulllineitems}
\phantomsection\label{\detokenize{api:src.hiveplot.HivePlotLayout.intra_axis_edges}}
\pysigstartsignatures
\pysigline
{\sphinxbfcode{\sphinxupquote{intra\_axis\_edges}}\sphinxbfcode{\sphinxupquote{\DUrole{p}{:}\DUrole{w}{ }list\DUrole{p}{{[}}tuple\DUrole{p}{{[}}int\DUrole{p}{,}\DUrole{w}{ }int\DUrole{p}{{]}}\DUrole{p}{{]}}}}}
\pysigstopsignatures
\end{fulllineitems}

\index{intra\_axis\_nodes (Attribut von src.hiveplot.HivePlotLayout)@\spxentry{intra\_axis\_nodes}\spxextra{Attribut von src.hiveplot.HivePlotLayout}}

\begin{fulllineitems}
\phantomsection\label{\detokenize{api:src.hiveplot.HivePlotLayout.intra_axis_nodes}}
\pysigstartsignatures
\pysigline
{\sphinxbfcode{\sphinxupquote{intra\_axis\_nodes}}\sphinxbfcode{\sphinxupquote{\DUrole{p}{:}\DUrole{w}{ }dict\DUrole{p}{{[}}int\DUrole{p}{,}\DUrole{w}{ }list\DUrole{p}{{[}}int\DUrole{p}{{]}}\DUrole{p}{{]}}}}}
\pysigstopsignatures
\end{fulllineitems}

\index{long\_edges (Attribut von src.hiveplot.HivePlotLayout)@\spxentry{long\_edges}\spxextra{Attribut von src.hiveplot.HivePlotLayout}}

\begin{fulllineitems}
\phantomsection\label{\detokenize{api:src.hiveplot.HivePlotLayout.long_edges}}
\pysigstartsignatures
\pysigline
{\sphinxbfcode{\sphinxupquote{long\_edges}}\sphinxbfcode{\sphinxupquote{\DUrole{p}{:}\DUrole{w}{ }set\DUrole{p}{{[}}tuple\DUrole{p}{{[}}int\DUrole{p}{,}\DUrole{w}{ }int\DUrole{p}{{]}}\DUrole{p}{{]}}}}}
\pysigstopsignatures
\end{fulllineitems}

\index{mixed\_nodes\_by\_axis (Attribut von src.hiveplot.HivePlotLayout)@\spxentry{mixed\_nodes\_by\_axis}\spxextra{Attribut von src.hiveplot.HivePlotLayout}}

\begin{fulllineitems}
\phantomsection\label{\detokenize{api:src.hiveplot.HivePlotLayout.mixed_nodes_by_axis}}
\pysigstartsignatures
\pysigline
{\sphinxbfcode{\sphinxupquote{mixed\_nodes\_by\_axis}}\sphinxbfcode{\sphinxupquote{\DUrole{p}{:}\DUrole{w}{ }dict\DUrole{p}{{[}}int\DUrole{p}{,}\DUrole{w}{ }list\DUrole{p}{{[}}int\DUrole{w}{ }\DUrole{p}{|}\DUrole{w}{ }str\DUrole{p}{{]}}\DUrole{p}{{]}}}}}
\pysigstopsignatures
\end{fulllineitems}

\index{name\_to\_id (Attribut von src.hiveplot.HivePlotLayout)@\spxentry{name\_to\_id}\spxextra{Attribut von src.hiveplot.HivePlotLayout}}

\begin{fulllineitems}
\phantomsection\label{\detokenize{api:src.hiveplot.HivePlotLayout.name_to_id}}
\pysigstartsignatures
\pysigline
{\sphinxbfcode{\sphinxupquote{name\_to\_id}}\sphinxbfcode{\sphinxupquote{\DUrole{p}{:}\DUrole{w}{ }dict\DUrole{p}{{[}}str\DUrole{p}{,}\DUrole{w}{ }int\DUrole{p}{{]}}}}}
\pysigstopsignatures
\end{fulllineitems}

\index{node\_groups (Attribut von src.hiveplot.HivePlotLayout)@\spxentry{node\_groups}\spxextra{Attribut von src.hiveplot.HivePlotLayout}}

\begin{fulllineitems}
\phantomsection\label{\detokenize{api:id2}}
\pysigstartsignatures
\pysigline
{\sphinxbfcode{\sphinxupquote{node\_groups}}\sphinxbfcode{\sphinxupquote{\DUrole{p}{:}\DUrole{w}{ }dict\DUrole{p}{{[}}int\DUrole{p}{,}\DUrole{w}{ }list\DUrole{p}{{[}}int\DUrole{w}{ }\DUrole{p}{|}\DUrole{w}{ }str\DUrole{p}{{]}}\DUrole{p}{{]}}}}}
\pysigstopsignatures
\end{fulllineitems}

\index{node\_groups\_dummies (Attribut von src.hiveplot.HivePlotLayout)@\spxentry{node\_groups\_dummies}\spxextra{Attribut von src.hiveplot.HivePlotLayout}}

\begin{fulllineitems}
\phantomsection\label{\detokenize{api:id3}}
\pysigstartsignatures
\pysigline
{\sphinxbfcode{\sphinxupquote{node\_groups\_dummies}}\sphinxbfcode{\sphinxupquote{\DUrole{p}{:}\DUrole{w}{ }dict\DUrole{p}{{[}}int\DUrole{p}{,}\DUrole{w}{ }list\DUrole{p}{{[}}str\DUrole{p}{{]}}\DUrole{p}{{]}}}}}
\pysigstopsignatures
\end{fulllineitems}

\index{node\_groups\_expanded (Attribut von src.hiveplot.HivePlotLayout)@\spxentry{node\_groups\_expanded}\spxextra{Attribut von src.hiveplot.HivePlotLayout}}

\begin{fulllineitems}
\phantomsection\label{\detokenize{api:src.hiveplot.HivePlotLayout.node_groups_expanded}}
\pysigstartsignatures
\pysigline
{\sphinxbfcode{\sphinxupquote{node\_groups\_expanded}}\sphinxbfcode{\sphinxupquote{\DUrole{p}{:}\DUrole{w}{ }dict\DUrole{p}{{[}}int\DUrole{p}{,}\DUrole{w}{ }list\DUrole{p}{{[}}int\DUrole{w}{ }\DUrole{p}{|}\DUrole{w}{ }str\DUrole{p}{{]}}\DUrole{p}{{]}}}}}
\pysigstopsignatures
\end{fulllineitems}

\index{num\_axes (Attribut von src.hiveplot.HivePlotLayout)@\spxentry{num\_axes}\spxextra{Attribut von src.hiveplot.HivePlotLayout}}

\begin{fulllineitems}
\phantomsection\label{\detokenize{api:id4}}
\pysigstartsignatures
\pysigline
{\sphinxbfcode{\sphinxupquote{num\_axes}}\sphinxbfcode{\sphinxupquote{\DUrole{p}{:}\DUrole{w}{ }int}}}
\pysigstopsignatures
\end{fulllineitems}

\index{post\_processing\_expansion() (Methode von src.hiveplot.HivePlotLayout)@\spxentry{post\_processing\_expansion()}\spxextra{Methode von src.hiveplot.HivePlotLayout}}

\begin{fulllineitems}
\phantomsection\label{\detokenize{api:src.hiveplot.HivePlotLayout.post_processing_expansion}}
\pysigstartsignatures
\pysiglinewithargsret
{\sphinxbfcode{\sphinxupquote{post\_processing\_expansion}}}
{\sphinxparam{\DUrole{n}{node\_axis\_map}}\sphinxparamcomma \sphinxparam{\DUrole{n}{dummy\_copy}\DUrole{o}{=}\DUrole{default_value}{None}}}
{}
\pysigstopsignatures
\sphinxAtStartPar
Die Funktion dient dem post\sphinxhyphen{}processing des Eingabegraphen nach der Pipeline, sodass bei diesem die Achsen expandiert werden. Dazu wird die Achse i (\textgreater{} 0) zu den Achsen mit KnotenIDs i und \sphinxhyphen{}i expandiert.  Die intra\sphinxhyphen{}axis Kanten werden symmetrisch zwischen den Achsenkopien gezeichnet.
Folgende Schritte werden durchgeführt:
1. Initialiseren der Dummysegmente, diese werden in der Funktion mutiert
2. Initialisieren der Map intra\_expandables mit Key AchsenID einer Kante und ihrer intra\sphinxhyphen{}axis Kanten (Achse wird nur aufgenommen, wenn es intra\sphinxhyphen{}axis Kanten gibt)
3. Initialisieren der Map inter\_expandables mit allen Kanten die genau einen Start oder Enknoten auf der expandierten Achse haben
4. Initielisieren der node\_groups\_expanded mit Keys ursprüngliche AchsenIDs und expandierte AchsenIDs und ihre in node\_groups enthaltenen Knotenlisten (expandierte Achse \sphinxhyphen{}i bekommt alle Knoten von i zugeordnet, jedoch werden die KnotenIDs negativ gesetzt, für virtuelle Knoten wird die neue Sequenznummer ermittelt und die neue Kante zwischen den virtuellen Knoten direkt in das Netwworkx Graphmodell hinzugefügt)
5. Update des Hiveplotlayouts mit neuer Achsenordnung und \sphinxhyphen{}anzahl und Update der edge\_axis\_map
6. Initialisieren einer axis\_position\_map mit Key AchsenID und Value Position der Achse in der neuen Achsenordnung
7. Behandlung der intra\sphinxhyphen{}axis Kanten
\begin{enumerate}
\sphinxsetlistlabels{\alph}{enumi}{enumii}{}{.}%
\item {} 
\sphinxAtStartPar
Entfernen der intra\sphinxhyphen{}axis Kanten aus dem Hiveplotlayout

\item {} 
\sphinxAtStartPar
Erstellen der neuen intra\sphinxhyphen{}axis Kanten zwischen den expandierten Achsen (Kante (u,v) auf Achse i wird symmetrisch gespiegelt zu den Kanten (\sphinxhyphen{}u, v) und (u, \sphinxhyphen{}v))

\item {} 
\sphinxAtStartPar
Einpflegen der neuen intra\sphinxhyphen{}axis Kanten in die Networkx Graphenstruktur des Hiveplotlayouts

\end{enumerate}
\begin{enumerate}
\sphinxsetlistlabels{\arabic}{enumi}{enumii}{}{.}%
\setcounter{enumi}{7}
\item {} \begin{description}
\sphinxlineitem{Behandlung der inter\sphinxhyphen{}axis Kanten}\begin{enumerate}
\sphinxsetlistlabels{\alph}{enumii}{enumiii}{}{.}%
\item {} 
\sphinxAtStartPar
Entfernen der inter\sphinxhyphen{}axis Kanten aus dem Hiveplotlayout und den Dummykantensegmenten

\item {} \begin{description}
\sphinxlineitem{Erstellen der neuen inter\sphinxhyphen{}axis Kanten, Fallbehandlung:}\begin{description}
\sphinxlineitem{I: Sonderfall: sowohl u als auch v auf expandierten achsen}\begin{description}
\sphinxlineitem{i.) v muss aktualisiert werden}
\sphinxAtStartPar
\textless{}1\textgreater{} v ist real und muss negiert werden
\textless{}2\textgreater{} v ist dummy und muss aktualisiert werden

\sphinxlineitem{ii.) u muss aktualisiert werden}
\sphinxAtStartPar
\textless{}1\textgreater{} u ist real und muss negiert werden
\textless{}2\textgreater{} u ist dummy und muss aktualisiert werden

\end{description}

\sphinxlineitem{II: u auf expandierter achse, v auf nicht expandierter achse}
\sphinxAtStartPar
i.) u ist real und muss negiert werden
ii.) u ist dummy und muss aktualisiert werden

\sphinxlineitem{III: v auf expandierter achse, u auf nicht expandierter achse}
\sphinxAtStartPar
i.) v ist real und muss negiert werden
ii.) v ist dummy und muss aktualisiert werden

\end{description}

\end{description}

\item {} 
\sphinxAtStartPar
Einpflegen der neuen inter\sphinxhyphen{}axis Kanten in die Networkx Graphenstruktur des Hiveplotlayouts

\end{enumerate}

\end{description}

\item {} 
\sphinxAtStartPar
alle neuen Knoten in hiveplot.graph.nodes speichern und subset zuordnen

\end{enumerate}
\begin{quote}\begin{description}
\sphinxlineitem{Parameter}\begin{itemize}
\item {} 
\sphinxAtStartPar
\sphinxstyleliteralstrong{\sphinxupquote{node\_axis\_map}} (\DUrole{sphinx_autodoc_typehints-type}{\sphinxcode{\sphinxupquote{dict}}{[}\sphinxcode{\sphinxupquote{int}} | \sphinxcode{\sphinxupquote{str}}, \sphinxcode{\sphinxupquote{int}}{]}}) \textendash{} Knoten\sphinxhyphen{}ID: Achsen\sphinxhyphen{}ID

\item {} 
\sphinxAtStartPar
\sphinxstyleliteralstrong{\sphinxupquote{dummy\_copy}} (\DUrole{sphinx_autodoc_typehints-type}{\sphinxcode{\sphinxupquote{list}}{[}\sphinxcode{\sphinxupquote{tuple}}{[}\sphinxcode{\sphinxupquote{int}} | \sphinxcode{\sphinxupquote{str}}, \sphinxcode{\sphinxupquote{int}} | \sphinxcode{\sphinxupquote{str}}{]}{]}}) \textendash{} eine Aktuelle Kopie der Dummysegmente, Default ist None

\end{itemize}

\sphinxlineitem{Rückgabetyp}
\sphinxAtStartPar
\DUrole{sphinx_autodoc_typehints-type}{\sphinxcode{\sphinxupquote{None}}}

\end{description}\end{quote}

\end{fulllineitems}

\index{pre\_processing\_expansion() (Methode von src.hiveplot.HivePlotLayout)@\spxentry{pre\_processing\_expansion()}\spxextra{Methode von src.hiveplot.HivePlotLayout}}

\begin{fulllineitems}
\phantomsection\label{\detokenize{api:src.hiveplot.HivePlotLayout.pre_processing_expansion}}
\pysigstartsignatures
\pysiglinewithargsret
{\sphinxbfcode{\sphinxupquote{pre\_processing\_expansion}}}
{\sphinxparam{\DUrole{n}{node\_axis\_map}}}
{}
\pysigstopsignatures
\sphinxAtStartPar
Die Funktion dient der Umsetzung unterschiedlicher Knotenordnungen auf expandierten Achsen. Die Funktion ist ausschließlich zum pre\sphinxhyphen{}processing des Eingabegraphen bevor etwaige Berechnungen der Pipeline durchgeführt werden gedacht. Dazu wird die Achse i zu den Achsen mit KnotenIDs i und \sphinxhyphen{}i expandiert.  Die intra\sphinxhyphen{}axis Kanten werden symmetrisch zwischen den Achsenkopien gezeichnet.
Folgende Schritte werden durchgeführt:
1. Initialisieren der Map intra\_expandables mit Key AchsenID einer Kante und ihrer intra\sphinxhyphen{}axis Kanten (Achse wird nur aufgenommen, wenn es intra\sphinxhyphen{}axis Kanten gibt)
2. Initialisieren der Map inter\_expandables mit allen Kanten die genau einen Start oder Enknoten auf der expandierten Achse haben
3. Initielisieren der node\_groups\_expanded mit Keys ursprüngliche AchsenIDs und expandierte AchsenIDs und ihre in node\_groups enthaltenen Knotenlisten (expandierte Achse \sphinxhyphen{}i bekommt alle Knoten von i zugeordnet, jedoch werden die KnotenIDs negativ gesetzt)
4. Update des Hiveplotlayouts mit neuer Achsenordnung und \sphinxhyphen{}anzahl
5. Initialisieren einer axis\_position\_map mit Key AchsenID und Value ist Position der Achse in der neuen Achsenordnung
6. Behandlung der intra\sphinxhyphen{}axis Kanten
\begin{enumerate}
\sphinxsetlistlabels{\alph}{enumi}{enumii}{}{.}%
\item {} 
\sphinxAtStartPar
Entfernen der intra\sphinxhyphen{}axis Kanten aus dem Hiveplotlayout

\item {} 
\sphinxAtStartPar
Erstellen der neuen intra\sphinxhyphen{}axis Kanten zwischen den expandierten Achsen (Kante (u,v) auf Achse i wird symmetrisch gespiegelt zu den Kanten (\sphinxhyphen{}u, v) und (u, \sphinxhyphen{}v))

\item {} 
\sphinxAtStartPar
Einpflegen der neuen intra\sphinxhyphen{}axis Kanten in die Networkx Graphenstruktur des Hiveplotlayouts

\end{enumerate}
\begin{enumerate}
\sphinxsetlistlabels{\arabic}{enumi}{enumii}{}{.}%
\setcounter{enumi}{6}
\item {} \begin{description}
\sphinxlineitem{Behandlung der inter\sphinxhyphen{}axis Kanten}\begin{enumerate}
\sphinxsetlistlabels{\alph}{enumii}{enumiii}{}{.}%
\item {} 
\sphinxAtStartPar
Entfernen der inter\sphinxhyphen{}axis Kanten aus dem Hiveplotlayout

\item {} \begin{description}
\sphinxlineitem{Erstellen der neuen inter\sphinxhyphen{}axis Kanten:}
\sphinxAtStartPar
I: Ziel\sphinxhyphen{} und Startknoten der Kante auf expandierter Achse: betrachte Kanten i/j mit Knoten u/v wenn i vor j in der Achsenordnung, Kante (u, v) zu (\sphinxhyphen{}u, v) andernfalls zu (u, \sphinxhyphen{}v) {[}i \sphinxhyphen{}\textgreater{} \sphinxhyphen{}i \sphinxhyphen{}\textgreater{} j \sphinxhyphen{}\textgreater{} \sphinxhyphen{}j{]}
II: Ziel\sphinxhyphen{} oder Startknoten der Kante auf expandierter Achse: betrachte Kanten i/j mit Knoten u/v (i expandiert), falls i vor j in der Achsenordnung (u,v) zu (\sphinxhyphen{}u, v), andernfalls Kante übernehmen {[}j \sphinxhyphen{}\textgreater{} i \sphinxhyphen{}\textgreater{} \sphinxhyphen{}i{]}

\end{description}

\item {} 
\sphinxAtStartPar
Einpflegen der neuen inter\sphinxhyphen{}axis Kanten in die Networkx Graphenstruktur des Hiveplotlayouts

\end{enumerate}

\end{description}

\item {} 
\sphinxAtStartPar
Einpflegen und Zuordnung zu Subsets der neuen Knoten in der Networkx Graphenstruktur des Hiveplotlayouts

\end{enumerate}
\begin{quote}\begin{description}
\sphinxlineitem{Parameter}
\sphinxAtStartPar
\sphinxstyleliteralstrong{\sphinxupquote{node\_axis\_map}} (\DUrole{sphinx_autodoc_typehints-type}{\sphinxcode{\sphinxupquote{dict}}{[}\sphinxcode{\sphinxupquote{int}} | \sphinxcode{\sphinxupquote{str}}, \sphinxcode{\sphinxupquote{int}}{]}}) \textendash{} Knoten\sphinxhyphen{}ID: Achsen\sphinxhyphen{}ID

\sphinxlineitem{Rückgabetyp}
\sphinxAtStartPar
\DUrole{sphinx_autodoc_typehints-type}{\sphinxcode{\sphinxupquote{None}}}

\end{description}\end{quote}

\end{fulllineitems}

\index{strict\_intra\_nodes\_by\_axis (Attribut von src.hiveplot.HivePlotLayout)@\spxentry{strict\_intra\_nodes\_by\_axis}\spxextra{Attribut von src.hiveplot.HivePlotLayout}}

\begin{fulllineitems}
\phantomsection\label{\detokenize{api:src.hiveplot.HivePlotLayout.strict_intra_nodes_by_axis}}
\pysigstartsignatures
\pysigline
{\sphinxbfcode{\sphinxupquote{strict\_intra\_nodes\_by\_axis}}\sphinxbfcode{\sphinxupquote{\DUrole{p}{:}\DUrole{w}{ }dict\DUrole{p}{{[}}int\DUrole{p}{,}\DUrole{w}{ }list\DUrole{p}{{[}}int\DUrole{w}{ }\DUrole{p}{|}\DUrole{w}{ }str\DUrole{p}{{]}}\DUrole{p}{{]}}}}}
\pysigstopsignatures
\end{fulllineitems}

\index{updated\_nodes() (Methode von src.hiveplot.HivePlotLayout)@\spxentry{updated\_nodes()}\spxextra{Methode von src.hiveplot.HivePlotLayout}}

\begin{fulllineitems}
\phantomsection\label{\detokenize{api:src.hiveplot.HivePlotLayout.updated_nodes}}
\pysigstartsignatures
\pysiglinewithargsret
{\sphinxbfcode{\sphinxupquote{updated\_nodes}}}
{\sphinxparam{\DUrole{n}{fused\_node\_groups}}}
{}
\pysigstopsignatures
\sphinxAtStartPar
Gibt die Knotenliste des Graphen zurück. Notwendig, weil die Knoten sich ändern und je nach Pipelineschritt die Knoten geupdated werden müssen.
\begin{quote}\begin{description}
\sphinxlineitem{Parameter}
\sphinxAtStartPar
\sphinxstyleliteralstrong{\sphinxupquote{fused\_node\_groups}} (\sphinxstyleliteralemphasis{\sphinxupquote{dict}}\sphinxstyleliteralemphasis{\sphinxupquote{{[}}}\sphinxstyleliteralemphasis{\sphinxupquote{int}}\sphinxstyleliteralemphasis{\sphinxupquote{, }}\sphinxstyleliteralemphasis{\sphinxupquote{list}}\sphinxstyleliteralemphasis{\sphinxupquote{{[}}}\sphinxstyleliteralemphasis{\sphinxupquote{int}}\sphinxstyleliteralemphasis{\sphinxupquote{ | }}\sphinxstyleliteralemphasis{\sphinxupquote{str}}\sphinxstyleliteralemphasis{\sphinxupquote{{]}}}\sphinxstyleliteralemphasis{\sphinxupquote{{]}}}) \textendash{} AchsenID: Knotenliste (real und virtuell)

\sphinxlineitem{Rückgabe}
\sphinxAtStartPar
Liste der realen und virtuellen Knoten

\sphinxlineitem{Rückgabetyp}
\sphinxAtStartPar
\DUrole{sphinx_autodoc_typehints-type}{\sphinxcode{\sphinxupquote{list}}{[}\sphinxcode{\sphinxupquote{int}}{]}}

\end{description}\end{quote}

\end{fulllineitems}


\end{fulllineitems}



\section{src.ip\_model module}
\label{\detokenize{api:module-src.ip_model}}\label{\detokenize{api:src-ip-model-module}}\index{module@\spxentry{module}!src.ip\_model@\spxentry{src.ip\_model}}\index{src.ip\_model@\spxentry{src.ip\_model}!module@\spxentry{module}}\index{delta\_mapping() (im Modul src.ip\_model)@\spxentry{delta\_mapping()}\spxextra{im Modul src.ip\_model}}

\begin{fulllineitems}
\phantomsection\label{\detokenize{api:src.ip_model.delta_mapping}}
\pysigstartsignatures
\pysiglinewithargsret
{\sphinxcode{\sphinxupquote{src.ip\_model.}}\sphinxbfcode{\sphinxupquote{delta\_mapping}}}
{\sphinxparam{\DUrole{n}{fused\_groups}}}
{}
\pysigstopsignatures
\sphinxAtStartPar
Dient der Initialisierung der Ordnungsvariablen (delta\textasciicircum{}i\_u,v) für das 1L2S\sphinxhyphen{}ILP. Erzeugt aus der Vereinigung von virtuellen und realen Knoten das Mapping.
\begin{quote}\begin{description}
\sphinxlineitem{Parameter}
\sphinxAtStartPar
\sphinxstyleliteralstrong{\sphinxupquote{layout}} ({\hyperref[\detokenize{api:src.hiveplot.HivePlotLayout}]{\sphinxcrossref{\sphinxstyleliteralemphasis{\sphinxupquote{HivePlotLayout}}}}}) \textendash{} das zugrundeliegende Hiveplotlayout

\sphinxlineitem{Rückgabe}
\sphinxAtStartPar
key: (u, v, alpha(u,v)), value: None

\sphinxlineitem{Rückgabetyp}
\sphinxAtStartPar
\DUrole{sphinx_autodoc_typehints-type}{\sphinxcode{\sphinxupquote{dict}}{[}\sphinxcode{\sphinxupquote{tuple}}{[}\sphinxcode{\sphinxupquote{int}} | \sphinxcode{\sphinxupquote{str}}, \sphinxcode{\sphinxupquote{...}}{]}, 1{]}}

\end{description}\end{quote}

\end{fulllineitems}

\index{get\_delta\_value() (im Modul src.ip\_model)@\spxentry{get\_delta\_value()}\spxextra{im Modul src.ip\_model}}

\begin{fulllineitems}
\phantomsection\label{\detokenize{api:src.ip_model.get_delta_value}}
\pysigstartsignatures
\pysiglinewithargsret
{\sphinxcode{\sphinxupquote{src.ip\_model.}}\sphinxbfcode{\sphinxupquote{get\_delta\_value}}}
{\sphinxparam{\DUrole{n}{delta}}\sphinxparamcomma \sphinxparam{\DUrole{n}{fixed\_delta}}\sphinxparamcomma \sphinxparam{\DUrole{n}{u}}\sphinxparamcomma \sphinxparam{\DUrole{n}{v}}\sphinxparamcomma \sphinxparam{\DUrole{n}{axis}}}
{}
\pysigstopsignatures
\sphinxAtStartPar
Die Funktion gibt das korrekte delta für fixierte oder nicht fixierte Ordnung zurück. Notwendig für 3b.
\begin{quote}\begin{description}
\sphinxlineitem{Rückgabetyp}
\sphinxAtStartPar
\DUrole{sphinx_autodoc_typehints-type}{\sphinxcode{\sphinxupquote{tuple}}{[}\sphinxcode{\sphinxupquote{int}} | \sphinxcode{\sphinxupquote{str}}, \sphinxcode{\sphinxupquote{int}} | \sphinxcode{\sphinxupquote{str}}, \sphinxcode{\sphinxupquote{int}}{]}}

\end{description}\end{quote}

\end{fulllineitems}

\index{ilp\_name() (im Modul src.ip\_model)@\spxentry{ilp\_name()}\spxextra{im Modul src.ip\_model}}

\begin{fulllineitems}
\phantomsection\label{\detokenize{api:src.ip_model.ilp_name}}
\pysigstartsignatures
\pysiglinewithargsret
{\sphinxcode{\sphinxupquote{src.ip\_model.}}\sphinxbfcode{\sphinxupquote{ilp\_name}}}
{\sphinxparam{\DUrole{o}{*}\DUrole{n}{parts}}}
{}
\pysigstopsignatures
\end{fulllineitems}

\index{induced\_crossings() (im Modul src.ip\_model)@\spxentry{induced\_crossings()}\spxextra{im Modul src.ip\_model}}

\begin{fulllineitems}
\phantomsection\label{\detokenize{api:src.ip_model.induced_crossings}}
\pysigstartsignatures
\pysiglinewithargsret
{\sphinxcode{\sphinxupquote{src.ip\_model.}}\sphinxbfcode{\sphinxupquote{induced\_crossings}}}
{\sphinxparam{\DUrole{n}{prob}}\sphinxparamcomma \sphinxparam{\DUrole{n}{delta\_static}}\sphinxparamcomma \sphinxparam{\DUrole{n}{fixed\_delta}}\sphinxparamcomma \sphinxparam{\DUrole{n}{sorted\_neighbors}}\sphinxparamcomma \sphinxparam{\DUrole{n}{pi\_fix}}\sphinxparamcomma \sphinxparam{\DUrole{n}{pi\_var\_id}}\sphinxparamcomma \sphinxparam{\DUrole{n}{pi\_var}}}
{}
\pysigstopsignatures
\sphinxAtStartPar
Ermittelt die Summanden der Zielfunktion für induzierte Kreuzungen. Für jedes Knotenpaar der variablen Achse werden alle Nachbarpaare auf
der festen Achse betrachtet. Kreuzungen werden durch die Linearisierung der Produkte der entsprechenden Ordnungsvariablen modelliert. Wir verwenden eine allgemeine
Linearisierung, da die Funktion get\_delta\_value sowohl feste Werte als auch PuLP\sphinxhyphen{}Variablen zurückgeben kann. Die Zielfunktion bleibt äquivalent zu Nöllenburg
und Wallinger, sie wird dadurch nur robust für den Solver.
\begin{quote}\begin{description}
\sphinxlineitem{Parameter}\begin{itemize}
\item {} 
\sphinxAtStartPar
\sphinxstyleliteralstrong{\sphinxupquote{prob}} (\DUrole{sphinx_autodoc_typehints-type}{\sphinxcode{\sphinxupquote{LpProblem}}}) \textendash{} das zugrunde liegende ILP\sphinxhyphen{}Modell

\item {} 
\sphinxAtStartPar
\sphinxstyleliteralstrong{\sphinxupquote{delta\_static}} (\DUrole{sphinx_autodoc_typehints-type}{\sphinxcode{\sphinxupquote{dict}}{[}\sphinxcode{\sphinxupquote{tuple}}{[}\sphinxcode{\sphinxupquote{int}} | \sphinxcode{\sphinxupquote{str}}, \sphinxcode{\sphinxupquote{int}} | \sphinxcode{\sphinxupquote{str}}, \sphinxcode{\sphinxupquote{int}}{]}, \sphinxcode{\sphinxupquote{object}}{]}}) \textendash{} bereits bekannte Ordnungsvariablen

\item {} 
\sphinxAtStartPar
\sphinxstyleliteralstrong{\sphinxupquote{fixed\_delta}} (\DUrole{sphinx_autodoc_typehints-type}{\sphinxcode{\sphinxupquote{dict}}{[}\sphinxcode{\sphinxupquote{tuple}}{[}\sphinxcode{\sphinxupquote{int}} | \sphinxcode{\sphinxupquote{str}}, \sphinxcode{\sphinxupquote{int}} | \sphinxcode{\sphinxupquote{str}}, \sphinxcode{\sphinxupquote{int}}{]}, \sphinxcode{\sphinxupquote{int}}{]}}) \textendash{} fixierte delta\sphinxhyphen{}Werte

\item {} 
\sphinxAtStartPar
\sphinxstyleliteralstrong{\sphinxupquote{sorted\_neighbors}} (\DUrole{sphinx_autodoc_typehints-type}{\sphinxcode{\sphinxupquote{dict}}{[}\sphinxcode{\sphinxupquote{tuple}}{[}\sphinxcode{\sphinxupquote{int}}, \sphinxcode{\sphinxupquote{int}} | \sphinxcode{\sphinxupquote{str}}{]}, \sphinxcode{\sphinxupquote{list}}{[}\sphinxcode{\sphinxupquote{int}} | \sphinxcode{\sphinxupquote{str}}{]}{]}}) \textendash{} Nachbarschaftsmap

\item {} 
\sphinxAtStartPar
\sphinxstyleliteralstrong{\sphinxupquote{pi\_fix}} (\DUrole{sphinx_autodoc_typehints-type}{\sphinxcode{\sphinxupquote{int}}}) \textendash{} ID feste Achse

\item {} 
\sphinxAtStartPar
\sphinxstyleliteralstrong{\sphinxupquote{pi\_var\_id}} (\DUrole{sphinx_autodoc_typehints-type}{\sphinxcode{\sphinxupquote{int}}}) \textendash{} ID variable Achse

\item {} 
\sphinxAtStartPar
\sphinxstyleliteralstrong{\sphinxupquote{pi\_var}} (\DUrole{sphinx_autodoc_typehints-type}{\sphinxcode{\sphinxupquote{list}}{[}\sphinxcode{\sphinxupquote{int}} | \sphinxcode{\sphinxupquote{str}}{]}}) \textendash{} Knoten der variablen Achse

\end{itemize}

\sphinxlineitem{Rückgabe}
\sphinxAtStartPar
lineare Zielfunktion zur Minimierung induzierter Kreuzungen

\sphinxlineitem{Rückgabetyp}
\sphinxAtStartPar
\DUrole{sphinx_autodoc_typehints-type}{\sphinxcode{\sphinxupquote{object}}}

\end{description}\end{quote}

\end{fulllineitems}

\index{ip\_model\_pipeline() (im Modul src.ip\_model)@\spxentry{ip\_model\_pipeline()}\spxextra{im Modul src.ip\_model}}

\begin{fulllineitems}
\phantomsection\label{\detokenize{api:src.ip_model.ip_model_pipeline}}
\pysigstartsignatures
\pysiglinewithargsret
{\sphinxcode{\sphinxupquote{src.ip\_model.}}\sphinxbfcode{\sphinxupquote{ip\_model\_pipeline}}}
{\sphinxparam{\DUrole{n}{layout}}\sphinxparamcomma \sphinxparam{\DUrole{n}{logger}}\sphinxparamcomma \sphinxparam{\DUrole{n}{threshold}\DUrole{o}{=}\DUrole{default_value}{10}}\sphinxparamcomma \sphinxparam{\DUrole{n}{paper\_like}\DUrole{o}{=}\DUrole{default_value}{True}}}
{}
\pysigstopsignatures
\sphinxAtStartPar
Die Funktion führt die gesamte Pipeline zur Kreuzungsminimierung durch, inklusive der Berechnung des 1L2S\sphinxhyphen{}ILP. Alle Schritte werden in\sphinxhyphen{}place am HivePlotLayout durchgeführt.
Ablauf paper\_like = false:
1. anfängliches Layout expandieren
2. isolierte Knoten entfernen
3. Segmentierung
4. 1L2S\sphinxhyphen{}ILP berechnen
5. Herstellen der Achsenornung und letztes Update der Datenstrukturen um die Visualisierung starten zu können.

\sphinxAtStartPar
Ablauf paper\_like = true:
1. isolierte Knoten entfernen
2. Segmentierung
3. 1L2S\sphinxhyphen{}ILP berechnen 3a
4. 1L2S\sphinxhyphen{}ILP berechnen 3b
5. Herstellen der Achsenornung und letztes Update der Datenstrukturen um die Visualisierung starten zu können.
\begin{quote}\begin{description}
\sphinxlineitem{Parameter}\begin{itemize}
\item {} 
\sphinxAtStartPar
\sphinxstyleliteralstrong{\sphinxupquote{layout}} (\DUrole{sphinx_autodoc_typehints-type}{\sphinxcode{\sphinxupquote{HivePlotLayout}}}) \textendash{} das zugrundeliegende Hiveplotlayout

\item {} 
\sphinxAtStartPar
\sphinxstyleliteralstrong{\sphinxupquote{threshold}} (\DUrole{sphinx_autodoc_typehints-type}{\sphinxcode{\sphinxupquote{int}}}) \textendash{} Anzahl der Sweeps bis zum Abbruch, Defaultwert ist 10.

\item {} 
\sphinxAtStartPar
\sphinxstyleliteralstrong{\sphinxupquote{layout\_expanded}} (\sphinxstyleliteralemphasis{\sphinxupquote{bool}}) \textendash{} dient der Unterscheidung, ob in der Pipeline mit expandierten Achsen gerechnet wird oder nicht, Default = False (nicht expandierter Fall)

\item {} 
\sphinxAtStartPar
\sphinxstyleliteralstrong{\sphinxupquote{paper\_like}} (\DUrole{sphinx_autodoc_typehints-type}{\sphinxcode{\sphinxupquote{bool}}}) \textendash{} dient der Unterscheidung zwischen identischer und unterschiedlicher Knotenordnung

\end{itemize}

\sphinxlineitem{Rückgabetyp}
\sphinxAtStartPar
\DUrole{sphinx_autodoc_typehints-type}{\sphinxcode{\sphinxupquote{None}}}

\end{description}\end{quote}

\end{fulllineitems}

\index{linearize\_product() (im Modul src.ip\_model)@\spxentry{linearize\_product()}\spxextra{im Modul src.ip\_model}}

\begin{fulllineitems}
\phantomsection\label{\detokenize{api:src.ip_model.linearize_product}}
\pysigstartsignatures
\pysiglinewithargsret
{\sphinxcode{\sphinxupquote{src.ip\_model.}}\sphinxbfcode{\sphinxupquote{linearize\_product}}}
{\sphinxparam{\DUrole{n}{prob}}\sphinxparamcomma \sphinxparam{\DUrole{n}{a}}\sphinxparamcomma \sphinxparam{\DUrole{n}{b}}\sphinxparamcomma \sphinxparam{\DUrole{n}{name}}}
{}
\pysigstopsignatures
\sphinxAtStartPar
Die Funktion kapselt die Linearisierung der quadratischen Terme in der Zielfunktion. Legt sowohl Variablen an als auch Nebenbedingungen fest.
\begin{quote}\begin{description}
\sphinxlineitem{Rückgabetyp}
\sphinxAtStartPar
\DUrole{sphinx_autodoc_typehints-type}{\sphinxcode{\sphinxupquote{LpVariable}}}

\end{description}\end{quote}

\end{fulllineitems}

\index{onelayer\_twosided\_optimization() (im Modul src.ip\_model)@\spxentry{onelayer\_twosided\_optimization()}\spxextra{im Modul src.ip\_model}}

\begin{fulllineitems}
\phantomsection\label{\detokenize{api:src.ip_model.onelayer_twosided_optimization}}
\pysigstartsignatures
\pysiglinewithargsret
{\sphinxcode{\sphinxupquote{src.ip\_model.}}\sphinxbfcode{\sphinxupquote{onelayer\_twosided\_optimization}}}
{\sphinxparam{\DUrole{n}{layout}}\sphinxparamcomma \sphinxparam{\DUrole{n}{neighborhood\_map}}\sphinxparamcomma \sphinxparam{\DUrole{n}{node\_axis\_map}}\sphinxparamcomma \sphinxparam{\DUrole{n}{threshold}\DUrole{o}{=}\DUrole{default_value}{10}}\sphinxparamcomma \sphinxparam{\DUrole{n}{layout\_expanded}\DUrole{o}{=}\DUrole{default_value}{False}}\sphinxparamcomma \sphinxparam{\DUrole{n}{paper\_like}\DUrole{o}{=}\DUrole{default_value}{False}}}
{}
\pysigstopsignatures
\sphinxAtStartPar
Enthält die gesamte Logik, um das One Layer Two Sided Integer Linear Program (1L2S ILP) zu definieren und zu berechnen. Die Funktion verändert das HiveplotLayout in\sphinxhyphen{}place. Die Pipeline besteht aus Sweeps.
Ein Sweep umfasst eine clockwise und eine counterclockwise Berechnung über jede Achse des Hiveplots. Folgende Schritte werden in der Funktion durchgeführt:
1. Initialisierung der Ordnungsvariablen (delta\textasciicircum{}i\_u,v) durch delta\_mapping()
2. Initialisierung der Sweepschleife, Abbruch bei Erreichen eines Thresholds
3. Danach schließt sich ein Sweep an erst in clockwise dann counterclockwise Richtung, beide habe das gleiche Schema:
\begin{enumerate}
\sphinxsetlistlabels{\alph}{enumi}{enumii}{}{.}%
\item {} 
\sphinxAtStartPar
Erstellen einer Probleminstanz und einer Kopie für die aktuelle Achse

\item {} \begin{description}
\sphinxlineitem{Initialisierung der Ordnungsvariablen für die aktuelle Achse (virtuelle und reale Knotenordnung werden in einem Schritt berechnet, aber die Achsenordnung real \textless{} virtuell muss berücksichtigt werden, da Gaps nicht explizit behandelt werden)}
\sphinxAtStartPar
I: Ordnungsvariablen getrennt für reale und virtuelle Achsenabschnitte
II: Ordnungsvariablen für real \textless{} virtuell = konstant 1
III: Ordnungsvariablen für virtuell \textless{} real = konstant 0

\end{description}

\item {} 
\sphinxAtStartPar
Erstellen der Zielfunktion speziell für die variable Achse

\item {} \begin{description}
\sphinxlineitem{Erstellen der Nebenbedingungen für:}
\sphinxAtStartPar
I: Antisymmetrie (im Paper implizit behandelt)
II: Transitivität

\end{description}

\item {} 
\sphinxAtStartPar
Lösen des ILP

\item {} 
\sphinxAtStartPar
Synchroniesieren der Lösung mit delta

\end{enumerate}
\begin{enumerate}
\sphinxsetlistlabels{\arabic}{enumi}{enumii}{}{.}%
\setcounter{enumi}{3}
\item {} 
\sphinxAtStartPar
Aktualisieren des Layouts durch Übersetzung von delta zurück in die Achsenordnung, da die variable Achse auf einem geupdateten Stand ermittelt werden muss, wird dies am Ende eines Sweeps durchgeführt

\end{enumerate}
\begin{quote}\begin{description}
\sphinxlineitem{Parameter}\begin{itemize}
\item {} 
\sphinxAtStartPar
\sphinxstyleliteralstrong{\sphinxupquote{layout}} (\DUrole{sphinx_autodoc_typehints-type}{\sphinxcode{\sphinxupquote{HivePlotLayout}}}) \textendash{} das zugrundeliegende HivePlotLayout

\item {} 
\sphinxAtStartPar
\sphinxstyleliteralstrong{\sphinxupquote{neighborhood\_map}} (\DUrole{sphinx_autodoc_typehints-type}{\sphinxcode{\sphinxupquote{dict}}{[}\sphinxcode{\sphinxupquote{str}} | \sphinxcode{\sphinxupquote{int}}, \sphinxcode{\sphinxupquote{list}}{[}\sphinxcode{\sphinxupquote{int}} | \sphinxcode{\sphinxupquote{str}}{]}{]}}) \textendash{} KnotenID: Liste von proper Nachbarn

\item {} 
\sphinxAtStartPar
\sphinxstyleliteralstrong{\sphinxupquote{node\_axis\_map}} (\DUrole{sphinx_autodoc_typehints-type}{\sphinxcode{\sphinxupquote{dict}}{[}\sphinxcode{\sphinxupquote{str}} | \sphinxcode{\sphinxupquote{int}}, \sphinxcode{\sphinxupquote{int}}{]}}) \textendash{} KnotenID: Achse

\item {} 
\sphinxAtStartPar
\sphinxstyleliteralstrong{\sphinxupquote{threshold}} (\DUrole{sphinx_autodoc_typehints-type}{\sphinxcode{\sphinxupquote{int}}}) \textendash{} Anzahl der Sweeps (=1x cw+ 1x ccw), Defaultwert ist 10

\item {} 
\sphinxAtStartPar
\sphinxstyleliteralstrong{\sphinxupquote{layout\_expanded}} (\DUrole{sphinx_autodoc_typehints-type}{\sphinxcode{\sphinxupquote{bool}}}) \textendash{} dient der Unterscheidung, ob in der Pipeline mit expandierten Achsen gerechnet wird oder nicht, Default = False (nicht expandierter Fall)

\end{itemize}

\sphinxlineitem{Rückgabetyp}
\sphinxAtStartPar
\DUrole{sphinx_autodoc_typehints-type}{\sphinxcode{\sphinxupquote{None}}}

\end{description}\end{quote}

\end{fulllineitems}

\index{onelayer\_twosided\_optimization\_3b() (im Modul src.ip\_model)@\spxentry{onelayer\_twosided\_optimization\_3b()}\spxextra{im Modul src.ip\_model}}

\begin{fulllineitems}
\phantomsection\label{\detokenize{api:src.ip_model.onelayer_twosided_optimization_3b}}
\pysigstartsignatures
\pysiglinewithargsret
{\sphinxcode{\sphinxupquote{src.ip\_model.}}\sphinxbfcode{\sphinxupquote{onelayer\_twosided\_optimization\_3b}}}
{\sphinxparam{\DUrole{n}{layout}}\sphinxparamcomma \sphinxparam{\DUrole{n}{neighborhood\_map}}\sphinxparamcomma \sphinxparam{\DUrole{n}{node\_axis\_map}}\sphinxparamcomma \sphinxparam{\DUrole{n}{threshold}\DUrole{o}{=}\DUrole{default_value}{10}}\sphinxparamcomma \sphinxparam{\DUrole{n}{layout\_expanded}\DUrole{o}{=}\DUrole{default_value}{False}}\sphinxparamcomma \sphinxparam{\DUrole{n}{paper\_like}\DUrole{o}{=}\DUrole{default_value}{False}}}
{}
\pysigstopsignatures
\sphinxAtStartPar
Die Funktion arbeitet analog zu onelayer\_twosided\_optimization(), berücksichtigt allerdings die zuvor fixierte globale Ordnung. Die Funktion wird explizit für Pipelineschritt 3b benötigt.
\begin{quote}\begin{description}
\sphinxlineitem{Parameter}\begin{itemize}
\item {} 
\sphinxAtStartPar
\sphinxstyleliteralstrong{\sphinxupquote{layout}} (\DUrole{sphinx_autodoc_typehints-type}{\sphinxcode{\sphinxupquote{HivePlotLayout}}}) \textendash{} das zugrundeliegende HivePlotLayout

\item {} 
\sphinxAtStartPar
\sphinxstyleliteralstrong{\sphinxupquote{neighborhood\_map}} (\DUrole{sphinx_autodoc_typehints-type}{\sphinxcode{\sphinxupquote{dict}}{[}\sphinxcode{\sphinxupquote{str}} | \sphinxcode{\sphinxupquote{int}}, \sphinxcode{\sphinxupquote{list}}{[}\sphinxcode{\sphinxupquote{int}} | \sphinxcode{\sphinxupquote{str}}{]}{]}}) \textendash{} KnotenID: Liste von proper Nachbarn

\item {} 
\sphinxAtStartPar
\sphinxstyleliteralstrong{\sphinxupquote{node\_axis\_map}} (\DUrole{sphinx_autodoc_typehints-type}{\sphinxcode{\sphinxupquote{dict}}{[}\sphinxcode{\sphinxupquote{str}} | \sphinxcode{\sphinxupquote{int}}, \sphinxcode{\sphinxupquote{int}}{]}}) \textendash{} KnotenID: Achse

\item {} 
\sphinxAtStartPar
\sphinxstyleliteralstrong{\sphinxupquote{threshold}} (\DUrole{sphinx_autodoc_typehints-type}{\sphinxcode{\sphinxupquote{int}}}) \textendash{} Anzahl der Sweeps (=1x cw+ 1x ccw), Defaultwert ist 10

\item {} 
\sphinxAtStartPar
\sphinxstyleliteralstrong{\sphinxupquote{layout\_expanded}} (\DUrole{sphinx_autodoc_typehints-type}{\sphinxcode{\sphinxupquote{bool}}}) \textendash{} dient der Unterscheidung, ob in der Pipeline mit expandierten Achsen gerechnet wird oder nicht, Default = False (nicht expandierter Fall)

\end{itemize}

\sphinxlineitem{Rückgabetyp}
\sphinxAtStartPar
\DUrole{sphinx_autodoc_typehints-type}{\sphinxcode{\sphinxupquote{None}}}

\end{description}\end{quote}

\end{fulllineitems}



\section{src.logger\_setup module}
\label{\detokenize{api:module-src.logger_setup}}\label{\detokenize{api:src-logger-setup-module}}\index{module@\spxentry{module}!src.logger\_setup@\spxentry{src.logger\_setup}}\index{src.logger\_setup@\spxentry{src.logger\_setup}!module@\spxentry{module}}\index{log() (im Modul src.logger\_setup)@\spxentry{log()}\spxextra{im Modul src.logger\_setup}}

\begin{fulllineitems}
\phantomsection\label{\detokenize{api:src.logger_setup.log}}
\pysigstartsignatures
\pysiglinewithargsret
{\sphinxcode{\sphinxupquote{src.logger\_setup.}}\sphinxbfcode{\sphinxupquote{log}}}
{\sphinxparam{\DUrole{n}{logger}}\sphinxparamcomma \sphinxparam{\DUrole{n}{msg}}}
{}
\pysigstopsignatures
\end{fulllineitems}

\index{setup\_logger() (im Modul src.logger\_setup)@\spxentry{setup\_logger()}\spxextra{im Modul src.logger\_setup}}

\begin{fulllineitems}
\phantomsection\label{\detokenize{api:src.logger_setup.setup_logger}}
\pysigstartsignatures
\pysiglinewithargsret
{\sphinxcode{\sphinxupquote{src.logger\_setup.}}\sphinxbfcode{\sphinxupquote{setup\_logger}}}
{\sphinxparam{\DUrole{n}{log\_dir}}\sphinxparamcomma \sphinxparam{\DUrole{n}{per\_session}\DUrole{o}{=}\DUrole{default_value}{True}}\sphinxparamcomma \sphinxparam{\DUrole{n}{file\_name}\DUrole{o}{=}\DUrole{default_value}{None}}}
{}
\pysigstopsignatures
\sphinxAtStartPar
Konfiguriert den Logger und initialisiert ihn. Kann für Batchläufe eine Datei schreiben und für Einzeldurchläufe mehrere.
\begin{quote}\begin{description}
\sphinxlineitem{Parameter}\begin{itemize}
\item {} 
\sphinxAtStartPar
\sphinxstyleliteralstrong{\sphinxupquote{log\_dir}} (\DUrole{sphinx_autodoc_typehints-type}{Path}) \textendash{} Ablagepfad

\item {} 
\sphinxAtStartPar
\sphinxstyleliteralstrong{\sphinxupquote{per\_session}} (\DUrole{sphinx_autodoc_typehints-type}{\sphinxcode{\sphinxupquote{bool}}}) \textendash{} Batch\sphinxhyphen{} oder Einzelauswahl. Defaults to True.

\item {} 
\sphinxAtStartPar
\sphinxstyleliteralstrong{\sphinxupquote{file\_name}} (\DUrole{sphinx_autodoc_typehints-type}{\sphinxcode{\sphinxupquote{str}} | \sphinxcode{\sphinxupquote{None}}}) \textendash{} Dateiname. Defaults to None.

\end{itemize}

\sphinxlineitem{Rückgabe}
\sphinxAtStartPar
der angelegte Log

\sphinxlineitem{Rückgabetyp}
\sphinxAtStartPar
\DUrole{sphinx_autodoc_typehints-type}{\sphinxcode{\sphinxupquote{Logger}}}

\end{description}\end{quote}

\end{fulllineitems}



\section{src.main module}
\label{\detokenize{api:module-src.main}}\label{\detokenize{api:src-main-module}}\index{module@\spxentry{module}!src.main@\spxentry{src.main}}\index{src.main@\spxentry{src.main}!module@\spxentry{module}}\index{init\_graph() (im Modul src.main)@\spxentry{init\_graph()}\spxextra{im Modul src.main}}

\begin{fulllineitems}
\phantomsection\label{\detokenize{api:src.main.init_graph}}
\pysigstartsignatures
\pysiglinewithargsret
{\sphinxcode{\sphinxupquote{src.main.}}\sphinxbfcode{\sphinxupquote{init\_graph}}}
{\sphinxparam{\DUrole{n}{original}}\sphinxparamcomma \sphinxparam{\DUrole{n}{name\_to\_id}}}
{}
\pysigstopsignatures
\end{fulllineitems}

\index{init\_hiveplot() (im Modul src.main)@\spxentry{init\_hiveplot()}\spxextra{im Modul src.main}}

\begin{fulllineitems}
\phantomsection\label{\detokenize{api:src.main.init_hiveplot}}
\pysigstartsignatures
\pysiglinewithargsret
{\sphinxcode{\sphinxupquote{src.main.}}\sphinxbfcode{\sphinxupquote{init\_hiveplot}}}
{\sphinxparam{\DUrole{n}{year}}\sphinxparamcomma \sphinxparam{\DUrole{n}{own\_pkl}}\sphinxparamcomma \sphinxparam{\DUrole{n}{partitions}}\sphinxparamcomma \sphinxparam{\DUrole{n}{logger}}}
{}
\pysigstopsignatures\begin{quote}\begin{description}
\sphinxlineitem{Rückgabetyp}
\sphinxAtStartPar
\DUrole{sphinx_autodoc_typehints-type}{\sphinxcode{\sphinxupquote{HivePlotLayout}}}

\end{description}\end{quote}

\end{fulllineitems}

\index{init\_logger() (im Modul src.main)@\spxentry{init\_logger()}\spxextra{im Modul src.main}}

\begin{fulllineitems}
\phantomsection\label{\detokenize{api:src.main.init_logger}}
\pysigstartsignatures
\pysiglinewithargsret
{\sphinxcode{\sphinxupquote{src.main.}}\sphinxbfcode{\sphinxupquote{init\_logger}}}
{\sphinxparam{\DUrole{n}{per\_session}}}
{}
\pysigstopsignatures
\end{fulllineitems}

\index{init\_original() (im Modul src.main)@\spxentry{init\_original()}\spxextra{im Modul src.main}}

\begin{fulllineitems}
\phantomsection\label{\detokenize{api:src.main.init_original}}
\pysigstartsignatures
\pysiglinewithargsret
{\sphinxcode{\sphinxupquote{src.main.}}\sphinxbfcode{\sphinxupquote{init\_original}}}
{\sphinxparam{\DUrole{n}{year}}\sphinxparamcomma \sphinxparam{\DUrole{n}{own\_pkl}\DUrole{o}{=}\DUrole{default_value}{None}}}
{}
\pysigstopsignatures
\sphinxAtStartPar
Lädt einen Graphen aus dem Cache oder einer eigenen Pickle\sphinxhyphen{}Datei.
\begin{quote}\begin{description}
\sphinxlineitem{Parameter}\begin{itemize}
\item {} 
\sphinxAtStartPar
\sphinxstyleliteralstrong{\sphinxupquote{year}} (\DUrole{sphinx_autodoc_typehints-type}{\sphinxcode{\sphinxupquote{int}}}) \textendash{} gewünschte GD\sphinxhyphen{}Jahr

\item {} 
\sphinxAtStartPar
\sphinxstyleliteralstrong{\sphinxupquote{own\_pkl}} (\DUrole{sphinx_autodoc_typehints-type}{\sphinxcode{\sphinxupquote{str}} | \sphinxcode{\sphinxupquote{None}}}) \textendash{} Dateiname einer .pkl\sphinxhyphen{}Datei in LOAD\_DIR. Muss ein nx.Graph\sphinxhyphen{}Objekt enthalten.

\end{itemize}

\sphinxlineitem{Rückgabe}
\sphinxAtStartPar
das Graph\sphinxhyphen{}Objekt

\sphinxlineitem{Rückgabetyp}
\sphinxAtStartPar
\DUrole{sphinx_autodoc_typehints-type}{\sphinxcode{\sphinxupquote{Graph}}}

\end{description}\end{quote}

\end{fulllineitems}

\index{main() (im Modul src.main)@\spxentry{main()}\spxextra{im Modul src.main}}

\begin{fulllineitems}
\phantomsection\label{\detokenize{api:src.main.main}}
\pysigstartsignatures
\pysiglinewithargsret
{\sphinxcode{\sphinxupquote{src.main.}}\sphinxbfcode{\sphinxupquote{main}}}
{}
{}
\pysigstopsignatures
\end{fulllineitems}

\index{pipeline() (im Modul src.main)@\spxentry{pipeline()}\spxextra{im Modul src.main}}

\begin{fulllineitems}
\phantomsection\label{\detokenize{api:src.main.pipeline}}
\pysigstartsignatures
\pysiglinewithargsret
{\sphinxcode{\sphinxupquote{src.main.}}\sphinxbfcode{\sphinxupquote{pipeline}}}
{\sphinxparam{\DUrole{n}{year}}\sphinxparamcomma \sphinxparam{\DUrole{n}{run}}\sphinxparamcomma \sphinxparam{\DUrole{n}{logger}}\sphinxparamcomma \sphinxparam{\DUrole{n}{output\_name}}\sphinxparamcomma \sphinxparam{\DUrole{n}{variant}}\sphinxparamcomma \sphinxparam{\DUrole{n}{paper\_like}}\sphinxparamcomma \sphinxparam{\DUrole{n}{partitions}\DUrole{o}{=}\DUrole{default_value}{8}}\sphinxparamcomma \sphinxparam{\DUrole{n}{threshold}\DUrole{o}{=}\DUrole{default_value}{5}}\sphinxparamcomma \sphinxparam{\DUrole{n}{save}\DUrole{o}{=}\DUrole{default_value}{False}}\sphinxparamcomma \sphinxparam{\DUrole{n}{debug}\DUrole{o}{=}\DUrole{default_value}{False}}\sphinxparamcomma \sphinxparam{\DUrole{n}{batch}\DUrole{o}{=}\DUrole{default_value}{False}}\sphinxparamcomma \sphinxparam{\DUrole{n}{own\_pkl}\DUrole{o}{=}\DUrole{default_value}{None}}\sphinxparamcomma \sphinxparam{\DUrole{n}{gui}\DUrole{o}{=}\DUrole{default_value}{False}}}
{}
\pysigstopsignatures
\end{fulllineitems}

\index{save\_pkl() (im Modul src.main)@\spxentry{save\_pkl()}\spxextra{im Modul src.main}}

\begin{fulllineitems}
\phantomsection\label{\detokenize{api:src.main.save_pkl}}
\pysigstartsignatures
\pysiglinewithargsret
{\sphinxcode{\sphinxupquote{src.main.}}\sphinxbfcode{\sphinxupquote{save\_pkl}}}
{\sphinxparam{\DUrole{n}{hiveplot}}\sphinxparamcomma \sphinxparam{\DUrole{n}{name}}\sphinxparamcomma \sphinxparam{\DUrole{n}{save}}\sphinxparamcomma \sphinxparam{\DUrole{n}{logger}}}
{}
\pysigstopsignatures\begin{quote}\begin{description}
\sphinxlineitem{Rückgabetyp}
\sphinxAtStartPar
\DUrole{sphinx_autodoc_typehints-type}{\sphinxcode{\sphinxupquote{None}}}

\end{description}\end{quote}

\end{fulllineitems}

\index{save\_rendered\_hiveplot() (im Modul src.main)@\spxentry{save\_rendered\_hiveplot()}\spxextra{im Modul src.main}}

\begin{fulllineitems}
\phantomsection\label{\detokenize{api:src.main.save_rendered_hiveplot}}
\pysigstartsignatures
\pysiglinewithargsret
{\sphinxcode{\sphinxupquote{src.main.}}\sphinxbfcode{\sphinxupquote{save\_rendered\_hiveplot}}}
{\sphinxparam{\DUrole{n}{svg\_path}}\sphinxparamcomma \sphinxparam{\DUrole{n}{year}}\sphinxparamcomma \sphinxparam{\DUrole{n}{logger}}}
{}
\pysigstopsignatures\begin{quote}\begin{description}
\sphinxlineitem{Rückgabetyp}
\sphinxAtStartPar
\DUrole{sphinx_autodoc_typehints-type}{\sphinxcode{\sphinxupquote{None}}}

\end{description}\end{quote}

\end{fulllineitems}

\index{step\_ordering() (im Modul src.main)@\spxentry{step\_ordering()}\spxextra{im Modul src.main}}

\begin{fulllineitems}
\phantomsection\label{\detokenize{api:src.main.step_ordering}}
\pysigstartsignatures
\pysiglinewithargsret
{\sphinxcode{\sphinxupquote{src.main.}}\sphinxbfcode{\sphinxupquote{step\_ordering}}}
{\sphinxparam{\DUrole{n}{hiveplot}}\sphinxparamcomma \sphinxparam{\DUrole{n}{logger}}}
{}
\pysigstopsignatures\begin{quote}\begin{description}
\sphinxlineitem{Rückgabetyp}
\sphinxAtStartPar
\DUrole{sphinx_autodoc_typehints-type}{\sphinxcode{\sphinxupquote{None}}}

\end{description}\end{quote}

\end{fulllineitems}



\section{src.ordering module}
\label{\detokenize{api:module-src.ordering}}\label{\detokenize{api:src-ordering-module}}\index{module@\spxentry{module}!src.ordering@\spxentry{src.ordering}}\index{src.ordering@\spxentry{src.ordering}!module@\spxentry{module}}\index{brute\_force\_ordering() (im Modul src.ordering)@\spxentry{brute\_force\_ordering()}\spxextra{im Modul src.ordering}}

\begin{fulllineitems}
\phantomsection\label{\detokenize{api:src.ordering.brute_force_ordering}}
\pysigstartsignatures
\pysiglinewithargsret
{\sphinxcode{\sphinxupquote{src.ordering.}}\sphinxbfcode{\sphinxupquote{brute\_force\_ordering}}}
{\sphinxparam{\DUrole{n}{ordering}}\sphinxparamcomma \sphinxparam{\DUrole{n}{node\_grps}}\sphinxparamcomma \sphinxparam{\DUrole{n}{edges}}}
{}
\pysigstopsignatures
\sphinxAtStartPar
Berechnet die optimale Knotenordnung nach der Kostenfunktion über alle Permutationen von k.
Diese Funktion wird für Vergleichs\sphinxhyphen{} und Testzwecke weiterhin bereitgestellt.
:param ordering: aktuelle Achsenordnung (phi)
:type ordering: \DUrole{sphinx_autodoc_typehints-type}{\sphinxcode{\sphinxupquote{list}}{[}\sphinxcode{\sphinxupquote{int}}{]}}
:param node\_grps: key: Achse, value: Knoten
:type node\_grps: \DUrole{sphinx_autodoc_typehints-type}{\sphinxcode{\sphinxupquote{dict}}}
:param edges: Knotenliste
:type edges: \DUrole{sphinx_autodoc_typehints-type}{\sphinxcode{\sphinxupquote{list}}}
\begin{quote}\begin{description}
\sphinxlineitem{Rückgabe}
\sphinxAtStartPar
Optimierte Achsenordnung

\sphinxlineitem{Rückgabetyp}
\sphinxAtStartPar
\DUrole{sphinx_autodoc_typehints-type}{\sphinxcode{\sphinxupquote{tuple}}{[}\sphinxcode{\sphinxupquote{int}}{]}}

\end{description}\end{quote}

\end{fulllineitems}

\index{ip\_ordering() (im Modul src.ordering)@\spxentry{ip\_ordering()}\spxextra{im Modul src.ordering}}

\begin{fulllineitems}
\phantomsection\label{\detokenize{api:src.ordering.ip_ordering}}
\pysigstartsignatures
\pysiglinewithargsret
{\sphinxcode{\sphinxupquote{src.ordering.}}\sphinxbfcode{\sphinxupquote{ip\_ordering}}}
{\sphinxparam{\DUrole{n}{layout}}}
{}
\pysigstopsignatures\begin{quote}\begin{description}
\sphinxlineitem{Rückgabetyp}
\sphinxAtStartPar
\DUrole{sphinx_autodoc_typehints-type}{\sphinxcode{\sphinxupquote{list}}{[}\sphinxcode{\sphinxupquote{int}}{]}}

\end{description}\end{quote}

\end{fulllineitems}

\index{node\_to\_axis() (im Modul src.ordering)@\spxentry{node\_to\_axis()}\spxextra{im Modul src.ordering}}

\begin{fulllineitems}
\phantomsection\label{\detokenize{api:src.ordering.node_to_axis}}
\pysigstartsignatures
\pysiglinewithargsret
{\sphinxcode{\sphinxupquote{src.ordering.}}\sphinxbfcode{\sphinxupquote{node\_to\_axis}}}
{\sphinxparam{\DUrole{n}{node}}\sphinxparamcomma \sphinxparam{\DUrole{n}{node\_grps}}\sphinxparamcomma \sphinxparam{\DUrole{n}{position}\DUrole{o}{=}\DUrole{default_value}{False}}}
{}
\pysigstopsignatures
\sphinxAtStartPar
Gibt an auf welcher Achse sich ein Knoten befindet. Funktion wird nur zum Initialisieren verwendet. Für Iterationen in der Pipeline ist sie nicht performant. Dafür werden Mappings benutzt (z.B. siehe ordering.node\_to\_axis\_maps)
\begin{quote}\begin{description}
\sphinxlineitem{Parameter}\begin{itemize}
\item {} 
\sphinxAtStartPar
\sphinxstyleliteralstrong{\sphinxupquote{node}} (\sphinxstyleliteralemphasis{\sphinxupquote{int}}) \textendash{} Knoten ID

\item {} 
\sphinxAtStartPar
\sphinxstyleliteralstrong{\sphinxupquote{node\_grps}} (\sphinxstyleliteralemphasis{\sphinxupquote{dict}}) \textendash{} key: Achse, value: Knoten

\item {} 
\sphinxAtStartPar
\sphinxstyleliteralstrong{\sphinxupquote{position}} (\sphinxstyleliteralemphasis{\sphinxupquote{bool}}) \textendash{} Rückgabe der Achse, true Rückgabe der Position der Achse in der Achsenordnung (phi). False ist default.

\end{itemize}

\sphinxlineitem{Verursacht}
\sphinxAtStartPar
\sphinxstyleliteralstrong{\sphinxupquote{ValueError}} \textendash{} Der Knoten kann keiner Achse zugeordnet werden.

\sphinxlineitem{Rückgabe}
\sphinxAtStartPar
Achsen ID

\sphinxlineitem{Rückgabetyp}
\sphinxAtStartPar
int

\end{description}\end{quote}

\end{fulllineitems}

\index{node\_to\_axis\_maps() (im Modul src.ordering)@\spxentry{node\_to\_axis\_maps()}\spxextra{im Modul src.ordering}}

\begin{fulllineitems}
\phantomsection\label{\detokenize{api:src.ordering.node_to_axis_maps}}
\pysigstartsignatures
\pysiglinewithargsret
{\sphinxcode{\sphinxupquote{src.ordering.}}\sphinxbfcode{\sphinxupquote{node\_to\_axis\_maps}}}
{\sphinxparam{\DUrole{n}{layout}}\sphinxparamcomma \sphinxparam{\DUrole{n}{fused\_node\_groups}}}
{}
\pysigstopsignatures
\sphinxAtStartPar
Die Funktion berechnet zwei Maps um für einen gegeben Knoten sehr schnell 1. alpha(Knoten) und 2. phi{[}alpha(Knoten){]} zu bestimmen.
\begin{quote}\begin{description}
\sphinxlineitem{Parameter}\begin{itemize}
\item {} 
\sphinxAtStartPar
\sphinxstyleliteralstrong{\sphinxupquote{layout}} (\DUrole{sphinx_autodoc_typehints-type}{\sphinxcode{\sphinxupquote{HivePlotLayout}}}) \textendash{} Das Hiveplotlayout

\item {} 
\sphinxAtStartPar
\sphinxstyleliteralstrong{\sphinxupquote{fused\_node\_groups}} (\DUrole{sphinx_autodoc_typehints-type}{\sphinxcode{\sphinxupquote{dict}}{[}\sphinxcode{\sphinxupquote{int}}, \sphinxcode{\sphinxupquote{list}}{[}\sphinxcode{\sphinxupquote{int}} | \sphinxcode{\sphinxupquote{str}}{]}{]}}) \textendash{} Achse: Knotenliste (virtuell und real)

\end{itemize}

\sphinxlineitem{Rückgabe}
\sphinxAtStartPar
Knoten: Achse \& Knoten: Achsenposition in Phi

\sphinxlineitem{Rückgabetyp}
\sphinxAtStartPar
\DUrole{sphinx_autodoc_typehints-type}{\sphinxcode{\sphinxupquote{tuple}}{[}\sphinxcode{\sphinxupquote{dict}}{[}\sphinxcode{\sphinxupquote{int}} | \sphinxcode{\sphinxupquote{str}}, \sphinxcode{\sphinxupquote{int}}{]}, \sphinxcode{\sphinxupquote{dict}}{[}\sphinxcode{\sphinxupquote{int}} | \sphinxcode{\sphinxupquote{str}}, \sphinxcode{\sphinxupquote{int}}{]}{]}}

\end{description}\end{quote}

\end{fulllineitems}

\index{reordered\_node\_groups() (im Modul src.ordering)@\spxentry{reordered\_node\_groups()}\spxextra{im Modul src.ordering}}

\begin{fulllineitems}
\phantomsection\label{\detokenize{api:src.ordering.reordered_node_groups}}
\pysigstartsignatures
\pysiglinewithargsret
{\sphinxcode{\sphinxupquote{src.ordering.}}\sphinxbfcode{\sphinxupquote{reordered\_node\_groups}}}
{\sphinxparam{\DUrole{n}{node\_grps}}\sphinxparamcomma \sphinxparam{\DUrole{n}{new\_order}}}
{}
\pysigstopsignatures
\sphinxAtStartPar
Schnelles Umordnen nach Optimierung der Achsenordnung.
\begin{quote}\begin{description}
\sphinxlineitem{Rückgabetyp}
\sphinxAtStartPar
\DUrole{sphinx_autodoc_typehints-type}{\sphinxcode{\sphinxupquote{dict}}{[}\sphinxcode{\sphinxupquote{int}}, \sphinxcode{\sphinxupquote{list}}{[}\sphinxcode{\sphinxupquote{int}}{]}{]}}

\end{description}\end{quote}

\end{fulllineitems}



\section{src.partitioning module}
\label{\detokenize{api:module-src.partitioning}}\label{\detokenize{api:src-partitioning-module}}\index{module@\spxentry{module}!src.partitioning@\spxentry{src.partitioning}}\index{src.partitioning@\spxentry{src.partitioning}!module@\spxentry{module}}\index{clauset\_newman\_moore\_communities() (im Modul src.partitioning)@\spxentry{clauset\_newman\_moore\_communities()}\spxextra{im Modul src.partitioning}}

\begin{fulllineitems}
\phantomsection\label{\detokenize{api:src.partitioning.clauset_newman_moore_communities}}
\pysigstartsignatures
\pysiglinewithargsret
{\sphinxcode{\sphinxupquote{src.partitioning.}}\sphinxbfcode{\sphinxupquote{clauset\_newman\_moore\_communities}}}
{\sphinxparam{\DUrole{n}{graph}}\sphinxparamcomma \sphinxparam{\DUrole{n}{threshold}\DUrole{o}{=}\DUrole{default_value}{0}}}
{}
\pysigstopsignatures
\sphinxAtStartPar
Berechnet die Communities eines Graphen per Clauset\sphinxhyphen{}Newman\sphinxhyphen{}Moore Algorithmus und wandelt die frozenset Liste in ein dict um.
Falls ein Threshold übergeben wird, so wird die Anzahl der Communities auf den Threshold gestaucht. So wird sichergestellt, dass man
eine Anzahl Communities (Achsen) erhält, die noch einen übersichtlichen Hiveplot erzeugen lassen.
\begin{quote}\begin{description}
\sphinxlineitem{Parameter}\begin{itemize}
\item {} 
\sphinxAtStartPar
\sphinxstyleliteralstrong{\sphinxupquote{graph}} (\sphinxstyleliteralemphasis{\sphinxupquote{nx.Graph}}) \textendash{} Der Eingabegraph für den die Communities berechnet werden sollen.

\item {} 
\sphinxAtStartPar
\sphinxstyleliteralstrong{\sphinxupquote{threshold}} (\DUrole{sphinx_autodoc_typehints-type}{\sphinxcode{\sphinxupquote{int}}}) \textendash{} maximale Anzahl Communities, bei Default (=0) wird die natürliche Anzahl von Communities ermittelt und zurückgegeben

\end{itemize}

\sphinxlineitem{Rückgabe}
\sphinxAtStartPar
Ein Dictionary mit den Communities als Werte und ihren IDs als Schlüssel.

\sphinxlineitem{Rückgabetyp}
\sphinxAtStartPar
\DUrole{sphinx_autodoc_typehints-type}{\sphinxcode{\sphinxupquote{dict}}{[}\sphinxcode{\sphinxupquote{int}}, \sphinxcode{\sphinxupquote{list}}{[}\sphinxcode{\sphinxupquote{int}}{]}{]}}

\end{description}\end{quote}

\end{fulllineitems}

\index{louvain\_community\_detection() (im Modul src.partitioning)@\spxentry{louvain\_community\_detection()}\spxextra{im Modul src.partitioning}}

\begin{fulllineitems}
\phantomsection\label{\detokenize{api:src.partitioning.louvain_community_detection}}
\pysigstartsignatures
\pysiglinewithargsret
{\sphinxcode{\sphinxupquote{src.partitioning.}}\sphinxbfcode{\sphinxupquote{louvain\_community\_detection}}}
{\sphinxparam{\DUrole{n}{graph}}}
{}
\pysigstopsignatures
\sphinxAtStartPar
Berechnet die Communities eines Graphen per Louvain Community Detection und wandelt die frozenset Liste in ein dict um.
Diese Funktion wird für Vergleichs\sphinxhyphen{} und Testzwecke weiterhin bereitgestellt.
:type graph: nx.Graph
:param graph: Der Eingabegraph, für den die Communities berechnet werden sollen.
:type graph: nx.Graph
\begin{quote}\begin{description}
\sphinxlineitem{Rückgabe}
\sphinxAtStartPar
Ein Dictionary mit den Communities als Werte und ihren IDs als Schlüssel.

\sphinxlineitem{Rückgabetyp}
\sphinxAtStartPar
dict{[}int, list{[}int{]}{]}

\end{description}\end{quote}

\end{fulllineitems}



\section{src.renderer module}
\label{\detokenize{api:module-src.renderer}}\label{\detokenize{api:src-renderer-module}}\index{module@\spxentry{module}!src.renderer@\spxentry{src.renderer}}\index{src.renderer@\spxentry{src.renderer}!module@\spxentry{module}}\index{draw\_basis() (im Modul src.renderer)@\spxentry{draw\_basis()}\spxextra{im Modul src.renderer}}

\begin{fulllineitems}
\phantomsection\label{\detokenize{api:src.renderer.draw_basis}}
\pysigstartsignatures
\pysiglinewithargsret
{\sphinxcode{\sphinxupquote{src.renderer.}}\sphinxbfcode{\sphinxupquote{draw\_basis}}}
{\sphinxparam{\DUrole{n}{layout}}\sphinxparamcomma \sphinxparam{\DUrole{n}{node\_groups}}\sphinxparamcomma \sphinxparam{\DUrole{n}{edges}}\sphinxparamcomma \sphinxparam{\DUrole{n}{degree}\DUrole{o}{=}\DUrole{default_value}{False}}\sphinxparamcomma \sphinxparam{\DUrole{n}{id\_to\_label\_map}\DUrole{o}{=}\DUrole{default_value}{None}}\sphinxparamcomma \sphinxparam{\DUrole{n}{unordered}\DUrole{o}{=}\DUrole{default_value}{False}}\sphinxparamcomma \sphinxparam{\DUrole{n}{intra\_edges}\DUrole{o}{=}\DUrole{default_value}{None}}\sphinxparamcomma \sphinxparam{\DUrole{n}{axes\_labels}\DUrole{o}{=}\DUrole{default_value}{True}}}
{}
\pysigstopsignatures
\sphinxAtStartPar
Bündelt die Darstellung von Kanten, Knoten, Labels und Achsen.
\begin{quote}\begin{description}
\sphinxlineitem{Rückgabetyp}
\sphinxAtStartPar
\DUrole{sphinx_autodoc_typehints-type}{\sphinxcode{\sphinxupquote{tuple}}{[}\sphinxcode{\sphinxupquote{list}}{[}\sphinxcode{\sphinxupquote{str}}{]}, \sphinxcode{\sphinxupquote{list}}{[}\sphinxcode{\sphinxupquote{str}}{]}, \sphinxcode{\sphinxupquote{list}}{[}\sphinxcode{\sphinxupquote{str}}{]}, \sphinxcode{\sphinxupquote{list}}{[}\sphinxcode{\sphinxupquote{str}}{]}{]}}

\end{description}\end{quote}

\end{fulllineitems}

\index{hiveplot\_renderer() (im Modul src.renderer)@\spxentry{hiveplot\_renderer()}\spxextra{im Modul src.renderer}}

\begin{fulllineitems}
\phantomsection\label{\detokenize{api:src.renderer.hiveplot_renderer}}
\pysigstartsignatures
\pysiglinewithargsret
{\sphinxcode{\sphinxupquote{src.renderer.}}\sphinxbfcode{\sphinxupquote{hiveplot\_renderer}}}
{\sphinxparam{\DUrole{n}{name}}\sphinxparamcomma \sphinxparam{\DUrole{n}{layout}}\sphinxparamcomma \sphinxparam{\DUrole{n}{debug\_dir}}\sphinxparamcomma \sphinxparam{\DUrole{n}{degree}\DUrole{o}{=}\DUrole{default_value}{True}}\sphinxparamcomma \sphinxparam{\DUrole{n}{layout\_expanded}\DUrole{o}{=}\DUrole{default_value}{False}}\sphinxparamcomma \sphinxparam{\DUrole{n}{intra}\DUrole{o}{=}\DUrole{default_value}{False}}\sphinxparamcomma \sphinxparam{\DUrole{n}{node\_labels}\DUrole{o}{=}\DUrole{default_value}{True}}\sphinxparamcomma \sphinxparam{\DUrole{n}{axes\_labels}\DUrole{o}{=}\DUrole{default_value}{True}}\sphinxparamcomma \sphinxparam{\DUrole{n}{unordered}\DUrole{o}{=}\DUrole{default_value}{False}}\sphinxparamcomma \sphinxparam{\DUrole{n}{paper\_like}\DUrole{o}{=}\DUrole{default_value}{False}}}
{}
\pysigstopsignatures
\sphinxAtStartPar
Pipeline die das Rendern des fertig berechneten Hiveplotlayouts in eine Scalable Vector Graphics (SVG) realisiert.
\begin{quote}\begin{description}
\sphinxlineitem{Rückgabetyp}
\sphinxAtStartPar
\DUrole{sphinx_autodoc_typehints-type}{\sphinxcode{\sphinxupquote{Path}}}

\end{description}\end{quote}

\end{fulllineitems}

\index{render\_svg() (im Modul src.renderer)@\spxentry{render\_svg()}\spxextra{im Modul src.renderer}}

\begin{fulllineitems}
\phantomsection\label{\detokenize{api:src.renderer.render_svg}}
\pysigstartsignatures
\pysiglinewithargsret
{\sphinxcode{\sphinxupquote{src.renderer.}}\sphinxbfcode{\sphinxupquote{render\_svg}}}
{\sphinxparam{\DUrole{n}{filename}}\sphinxparamcomma \sphinxparam{\DUrole{n}{width}}\sphinxparamcomma \sphinxparam{\DUrole{n}{height}}\sphinxparamcomma \sphinxparam{\DUrole{n}{elements}}}
{}
\pysigstopsignatures
\sphinxAtStartPar
Setzt den Header der .svg Datei und hängt die Graphenelemente an und setzt die Schlussklausel. Datei wird neu erstellt bzw. überschrieben, falls sie bereits existiert.
\begin{quote}\begin{description}
\sphinxlineitem{Rückgabetyp}
\sphinxAtStartPar
\DUrole{sphinx_autodoc_typehints-type}{\sphinxcode{\sphinxupquote{None}}}

\end{description}\end{quote}

\end{fulllineitems}

\index{translate\_polar\_to\_carthesian() (im Modul src.renderer)@\spxentry{translate\_polar\_to\_carthesian()}\spxextra{im Modul src.renderer}}

\begin{fulllineitems}
\phantomsection\label{\detokenize{api:src.renderer.translate_polar_to_carthesian}}
\pysigstartsignatures
\pysiglinewithargsret
{\sphinxcode{\sphinxupquote{src.renderer.}}\sphinxbfcode{\sphinxupquote{translate\_polar\_to\_carthesian}}}
{\sphinxparam{\DUrole{n}{radius}}\sphinxparamcomma \sphinxparam{\DUrole{n}{angle}}\sphinxparamcomma \sphinxparam{\DUrole{n}{center\_x}\DUrole{o}{=}\DUrole{default_value}{500.0}}\sphinxparamcomma \sphinxparam{\DUrole{n}{center\_y}\DUrole{o}{=}\DUrole{default_value}{500.0}}}
{}
\pysigstopsignatures
\sphinxAtStartPar
Umrechnung der Polarkoordinate in ein karthesisches Tupel.
\begin{quote}\begin{description}
\sphinxlineitem{Rückgabetyp}
\sphinxAtStartPar
\DUrole{sphinx_autodoc_typehints-type}{\sphinxcode{\sphinxupquote{tuple}}{[}\sphinxcode{\sphinxupquote{float}}, \sphinxcode{\sphinxupquote{float}}{]}}

\end{description}\end{quote}

\end{fulllineitems}



\section{Module contents}
\label{\detokenize{api:module-src}}\label{\detokenize{api:module-contents}}\index{module@\spxentry{module}!src@\spxentry{src}}\index{src@\spxentry{src}!module@\spxentry{module}}
\sphinxAtStartPar
Top\sphinxhyphen{}level package for project source files.


\renewcommand{\indexname}{Python-Modulindex}
\begin{sphinxtheindex}
\let\bigletter\sphinxstyleindexlettergroup
\bigletter{s}
\item\relax\sphinxstyleindexentry{src}\sphinxstyleindexpageref{api:\detokenize{module-src}}
\item\relax\sphinxstyleindexentry{src.cost}\sphinxstyleindexpageref{api:\detokenize{module-src.cost}}
\item\relax\sphinxstyleindexentry{src.crossing\_minimization}\sphinxstyleindexpageref{api:\detokenize{module-src.crossing_minimization}}
\item\relax\sphinxstyleindexentry{src.dblp\_parser}\sphinxstyleindexpageref{api:\detokenize{module-src.dblp_parser}}
\item\relax\sphinxstyleindexentry{src.experiments}\sphinxstyleindexpageref{api:\detokenize{module-src.experiments}}
\item\relax\sphinxstyleindexentry{src.hiveplot}\sphinxstyleindexpageref{api:\detokenize{module-src.hiveplot}}
\item\relax\sphinxstyleindexentry{src.ip\_model}\sphinxstyleindexpageref{api:\detokenize{module-src.ip_model}}
\item\relax\sphinxstyleindexentry{src.logger\_setup}\sphinxstyleindexpageref{api:\detokenize{module-src.logger_setup}}
\item\relax\sphinxstyleindexentry{src.main}\sphinxstyleindexpageref{api:\detokenize{module-src.main}}
\item\relax\sphinxstyleindexentry{src.ordering}\sphinxstyleindexpageref{api:\detokenize{module-src.ordering}}
\item\relax\sphinxstyleindexentry{src.partitioning}\sphinxstyleindexpageref{api:\detokenize{module-src.partitioning}}
\item\relax\sphinxstyleindexentry{src.renderer}\sphinxstyleindexpageref{api:\detokenize{module-src.renderer}}
\end{sphinxtheindex}

\renewcommand{\indexname}{Stichwortverzeichnis}
\printindex
\end{document}